% Options for packages loaded elsewhere
\PassOptionsToPackage{unicode}{hyperref}
\PassOptionsToPackage{hyphens}{url}
\documentclass[
]{article}
\usepackage{xcolor}
\usepackage[margin=1in]{geometry}
\usepackage{amsmath,amssymb}
\setcounter{secnumdepth}{-\maxdimen} % remove section numbering
\usepackage{iftex}
\ifPDFTeX
  \usepackage[T1]{fontenc}
  \usepackage[utf8]{inputenc}
  \usepackage{textcomp} % provide euro and other symbols
\else % if luatex or xetex
  \usepackage{unicode-math} % this also loads fontspec
  \defaultfontfeatures{Scale=MatchLowercase}
  \defaultfontfeatures[\rmfamily]{Ligatures=TeX,Scale=1}
\fi
\usepackage{lmodern}
\ifPDFTeX\else
  % xetex/luatex font selection
\fi
% Use upquote if available, for straight quotes in verbatim environments
\IfFileExists{upquote.sty}{\usepackage{upquote}}{}
\IfFileExists{microtype.sty}{% use microtype if available
  \usepackage[]{microtype}
  \UseMicrotypeSet[protrusion]{basicmath} % disable protrusion for tt fonts
}{}
\makeatletter
\@ifundefined{KOMAClassName}{% if non-KOMA class
  \IfFileExists{parskip.sty}{%
    \usepackage{parskip}
  }{% else
    \setlength{\parindent}{0pt}
    \setlength{\parskip}{6pt plus 2pt minus 1pt}}
}{% if KOMA class
  \KOMAoptions{parskip=half}}
\makeatother
\usepackage{color}
\usepackage{fancyvrb}
\newcommand{\VerbBar}{|}
\newcommand{\VERB}{\Verb[commandchars=\\\{\}]}
\DefineVerbatimEnvironment{Highlighting}{Verbatim}{commandchars=\\\{\}}
% Add ',fontsize=\small' for more characters per line
\usepackage{framed}
\definecolor{shadecolor}{RGB}{248,248,248}
\newenvironment{Shaded}{\begin{snugshade}}{\end{snugshade}}
\newcommand{\AlertTok}[1]{\textcolor[rgb]{0.94,0.16,0.16}{#1}}
\newcommand{\AnnotationTok}[1]{\textcolor[rgb]{0.56,0.35,0.01}{\textbf{\textit{#1}}}}
\newcommand{\AttributeTok}[1]{\textcolor[rgb]{0.13,0.29,0.53}{#1}}
\newcommand{\BaseNTok}[1]{\textcolor[rgb]{0.00,0.00,0.81}{#1}}
\newcommand{\BuiltInTok}[1]{#1}
\newcommand{\CharTok}[1]{\textcolor[rgb]{0.31,0.60,0.02}{#1}}
\newcommand{\CommentTok}[1]{\textcolor[rgb]{0.56,0.35,0.01}{\textit{#1}}}
\newcommand{\CommentVarTok}[1]{\textcolor[rgb]{0.56,0.35,0.01}{\textbf{\textit{#1}}}}
\newcommand{\ConstantTok}[1]{\textcolor[rgb]{0.56,0.35,0.01}{#1}}
\newcommand{\ControlFlowTok}[1]{\textcolor[rgb]{0.13,0.29,0.53}{\textbf{#1}}}
\newcommand{\DataTypeTok}[1]{\textcolor[rgb]{0.13,0.29,0.53}{#1}}
\newcommand{\DecValTok}[1]{\textcolor[rgb]{0.00,0.00,0.81}{#1}}
\newcommand{\DocumentationTok}[1]{\textcolor[rgb]{0.56,0.35,0.01}{\textbf{\textit{#1}}}}
\newcommand{\ErrorTok}[1]{\textcolor[rgb]{0.64,0.00,0.00}{\textbf{#1}}}
\newcommand{\ExtensionTok}[1]{#1}
\newcommand{\FloatTok}[1]{\textcolor[rgb]{0.00,0.00,0.81}{#1}}
\newcommand{\FunctionTok}[1]{\textcolor[rgb]{0.13,0.29,0.53}{\textbf{#1}}}
\newcommand{\ImportTok}[1]{#1}
\newcommand{\InformationTok}[1]{\textcolor[rgb]{0.56,0.35,0.01}{\textbf{\textit{#1}}}}
\newcommand{\KeywordTok}[1]{\textcolor[rgb]{0.13,0.29,0.53}{\textbf{#1}}}
\newcommand{\NormalTok}[1]{#1}
\newcommand{\OperatorTok}[1]{\textcolor[rgb]{0.81,0.36,0.00}{\textbf{#1}}}
\newcommand{\OtherTok}[1]{\textcolor[rgb]{0.56,0.35,0.01}{#1}}
\newcommand{\PreprocessorTok}[1]{\textcolor[rgb]{0.56,0.35,0.01}{\textit{#1}}}
\newcommand{\RegionMarkerTok}[1]{#1}
\newcommand{\SpecialCharTok}[1]{\textcolor[rgb]{0.81,0.36,0.00}{\textbf{#1}}}
\newcommand{\SpecialStringTok}[1]{\textcolor[rgb]{0.31,0.60,0.02}{#1}}
\newcommand{\StringTok}[1]{\textcolor[rgb]{0.31,0.60,0.02}{#1}}
\newcommand{\VariableTok}[1]{\textcolor[rgb]{0.00,0.00,0.00}{#1}}
\newcommand{\VerbatimStringTok}[1]{\textcolor[rgb]{0.31,0.60,0.02}{#1}}
\newcommand{\WarningTok}[1]{\textcolor[rgb]{0.56,0.35,0.01}{\textbf{\textit{#1}}}}
\usepackage{longtable,booktabs,array}
\newcounter{none} % for unnumbered tables
\usepackage{calc} % for calculating minipage widths
% Correct order of tables after \paragraph or \subparagraph
\usepackage{etoolbox}
\makeatletter
\patchcmd\longtable{\par}{\if@noskipsec\mbox{}\fi\par}{}{}
\makeatother
% Allow footnotes in longtable head/foot
\IfFileExists{footnotehyper.sty}{\usepackage{footnotehyper}}{\usepackage{footnote}}
\makesavenoteenv{longtable}
\usepackage{graphicx}
\makeatletter
\newsavebox\pandoc@box
\newcommand*\pandocbounded[1]{% scales image to fit in text height/width
  \sbox\pandoc@box{#1}%
  \Gscale@div\@tempa{\textheight}{\dimexpr\ht\pandoc@box+\dp\pandoc@box\relax}%
  \Gscale@div\@tempb{\linewidth}{\wd\pandoc@box}%
  \ifdim\@tempb\p@<\@tempa\p@\let\@tempa\@tempb\fi% select the smaller of both
  \ifdim\@tempa\p@<\p@\scalebox{\@tempa}{\usebox\pandoc@box}%
  \else\usebox{\pandoc@box}%
  \fi%
}
% Set default figure placement to htbp
\def\fps@figure{htbp}
\makeatother
\setlength{\emergencystretch}{3em} % prevent overfull lines
\providecommand{\tightlist}{%
  \setlength{\itemsep}{0pt}\setlength{\parskip}{0pt}}
\usepackage{booktabs}
\usepackage{longtable}
\usepackage{array}
\usepackage{multirow}
\usepackage{wrapfig}
\usepackage{float}
\usepackage{colortbl}
\usepackage{pdflscape}
\usepackage{tabu}
\usepackage{threeparttable}
\usepackage{threeparttablex}
\usepackage[normalem]{ulem}
\usepackage{makecell}
\usepackage{xcolor}
\usepackage{bookmark}
\IfFileExists{xurl.sty}{\usepackage{xurl}}{} % add URL line breaks if available
\urlstyle{same}
\hypersetup{
  pdftitle={Module 03: Measurement --- Exploratory and Confirmatory Factor Analysis},
  hidelinks,
  pdfcreator={LaTeX via pandoc}}

\title{Module 03: Measurement --- Exploratory and Confirmatory Factor
Analysis}
\usepackage{etoolbox}
\makeatletter
\providecommand{\subtitle}[1]{% add subtitle to \maketitle
  \apptocmd{\@title}{\par {\large #1 \par}}{}{}
}
\makeatother
\subtitle{Computational Social Science --- Graduate \& Undergraduate}
\author{}
\date{\vspace{-2.5em}2026-04-04}

\begin{document}
\maketitle

{
\setcounter{tocdepth}{2}
\tableofcontents
}
\begin{center}\rule{0.5\linewidth}{0.5pt}\end{center}

\section{Overview: Why Measurement
Matters}\label{overview-why-measurement-matters}

Before you can study relationships between constructs --- do political
values predict civic engagement? does stress cause health decline? ---
you need to establish that you are actually measuring those constructs.
This is the measurement problem, and it is harder than it looks.

The things social scientists most care about --- personality,
intelligence, ideology, wellbeing, social capital, cultural capital ---
are \textbf{latent variables}: theoretical constructs that cannot be
directly observed but must be inferred from their effects on observable
indicators. You cannot see ``neuroticism'' directly. You can see that a
person worries frequently, gets irritated easily, has frequent mood
swings, and often feels blue. Those four responses, taken together, give
you information about an underlying tendency you call neuroticism.

Factor analysis is the primary tool for modelling latent variables. It
exists in two forms:

\begin{itemize}
\tightlist
\item
  \textbf{EFA (Exploratory Factor Analysis):} Used when you do not know
  (or want to discover) the factor structure. You let the data tell you
  how many factors there are and which items belong to which.
\item
  \textbf{CFA (Confirmatory Factor Analysis):} Used when you have a
  theoretically specified factor structure and want to test whether it
  fits the data.
\end{itemize}

\textbf{Learning goals:}

\begin{itemize}
\tightlist
\item
  Explain what a latent variable is and why averaging items is not
  equivalent
\item
  Check whether data is suitable for factor analysis (KMO, Bartlett)
\item
  Run EFA, decide on the number of factors, choose a rotation, and read
  every element of the output
\item
  Specify and fit a CFA in lavaan
\item
  Read and interpret the full lavaan CFA output, section by section
\item
  Evaluate model fit using the standard indices with their benchmarks
\item
  Compute reliability with Cronbach's alpha and McDonald's omega
\end{itemize}

\begin{center}\rule{0.5\linewidth}{0.5pt}\end{center}

\section{The Measurement Problem}\label{the-measurement-problem}

\subsection{Why You Cannot Just Average
Items}\label{why-you-cannot-just-average-items}

Suppose you want to measure ``social trust'' using three survey items:

\begin{enumerate}
\def\labelenumi{\arabic{enumi}.}
\tightlist
\item
  ``Most people can be trusted.'' (1--5 agree)
\item
  ``People try to be helpful.'' (1--5 agree)
\item
  ``Most people are fair.'' (1--5 agree)
\end{enumerate}

A common approach is to sum or average these:
\texttt{trust\_score\ =\ (item1\ +\ item2\ +\ item3)\ /\ 3}. This is
called a \textbf{composite score}. It has several problems:

\begin{itemize}
\tightlist
\item
  It treats all three items as equally informative --- but some items
  may be better indicators of the underlying construct than others.
\item
  It assumes all measurement error is random --- but items can have
  systematic biases (e.g., item 1 is susceptible to mood effects not
  related to trust).
\item
  It gives no information about how reliably the composite measures the
  construct.
\item
  It cannot tell you whether the three items are actually measuring the
  same underlying thing.
\end{itemize}

Factor analysis solves all of these problems. Instead of assuming
equality, it \emph{estimates} the relationship between each item and the
underlying factor. Instead of ignoring error, it explicitly models it.
Instead of assuming a single factor, it tests whether one factor fits
the data.

\subsection{The Latent Variable Model}\label{the-latent-variable-model}

The single-factor model says: \[y_i = \lambda_i \eta + \varepsilon_i\]

\begin{itemize}
\tightlist
\item
  \(y_i\): the observed score on item \(i\)
\item
  \(\eta\): the latent factor (unobserved; what you actually want to
  measure)
\item
  \(\lambda_i\): the \textbf{factor loading} --- how strongly item \(i\)
  is related to \(\eta\)
\item
  \(\varepsilon_i\): \textbf{unique variance} --- the part of item \(i\)
  not explained by \(\eta\) (measurement error + item-specific content)
\end{itemize}

The \textbf{communality} of item \(i\) is \(h_i^2 = \lambda_i^2\) (in
standardised form) --- the proportion of item variance explained by the
factor. The \textbf{uniqueness} is \(u_i^2 = 1 - h_i^2\) --- the
proportion \emph{not} explained.

\textbf{Local independence} is the key assumption: within a given level
of \(\eta\), the items are uncorrelated. All item intercorrelations
arise from the shared latent factor. If two items remain correlated
after controlling for the factor, the model is misspecified.

\begin{center}\rule{0.5\linewidth}{0.5pt}\end{center}

\section{The IPIP Big Five Dataset}\label{the-ipip-big-five-dataset}

We use \texttt{bfi} from the \texttt{psych} package: 2,800 participants
× 25 personality items + demographics. Items measure the Big Five
personality traits (OCEAN: Openness, Conscientiousness, Extraversion,
Agreeableness, Neuroticism), 5 items each, on a 6-point scale (1 = Very
Inaccurate, 6 = Very Accurate).

\begin{Shaded}
\begin{Highlighting}[]
\FunctionTok{library}\NormalTok{(psych)}
\FunctionTok{data}\NormalTok{(bfi)}

\FunctionTok{cat}\NormalTok{(}\StringTok{"Dataset dimensions:"}\NormalTok{, }\FunctionTok{nrow}\NormalTok{(bfi), }\StringTok{"participants ×"}\NormalTok{, }\FunctionTok{ncol}\NormalTok{(bfi), }\StringTok{"variables}\SpecialCharTok{\textbackslash{}n}\StringTok{"}\NormalTok{)}
\end{Highlighting}
\end{Shaded}

\begin{verbatim}
## Dataset dimensions: 2800 participants × 28 variables
\end{verbatim}

\begin{Shaded}
\begin{Highlighting}[]
\FunctionTok{cat}\NormalTok{(}\StringTok{"First 25 columns: personality items (A1–A5, C1–C5, E1–E5, N1–N5, O1–O5)}\SpecialCharTok{\textbackslash{}n}\StringTok{"}\NormalTok{)}
\end{Highlighting}
\end{Shaded}

\begin{verbatim}
## First 25 columns: personality items (A1–A5, C1–C5, E1–E5, N1–N5, O1–O5)
\end{verbatim}

\begin{Shaded}
\begin{Highlighting}[]
\FunctionTok{cat}\NormalTok{(}\StringTok{"Remaining: gender, education, age}\SpecialCharTok{\textbackslash{}n\textbackslash{}n}\StringTok{"}\NormalTok{)}
\end{Highlighting}
\end{Shaded}

\begin{verbatim}
## Remaining: gender, education, age
\end{verbatim}

\begin{Shaded}
\begin{Highlighting}[]
\CommentTok{\# A quick look at one item to understand the scale}
\FunctionTok{cat}\NormalTok{(}\StringTok{"Distribution of item E3 (\textquotesingle{}Know how to captivate people\textquotesingle{}):}\SpecialCharTok{\textbackslash{}n}\StringTok{"}\NormalTok{)}
\end{Highlighting}
\end{Shaded}

\begin{verbatim}
## Distribution of item E3 ('Know how to captivate people'):
\end{verbatim}

\begin{Shaded}
\begin{Highlighting}[]
\FunctionTok{table}\NormalTok{(bfi}\SpecialCharTok{$}\NormalTok{E3, }\AttributeTok{useNA=}\StringTok{"ifany"}\NormalTok{) }\SpecialCharTok{|\textgreater{}} \FunctionTok{print}\NormalTok{()}
\end{Highlighting}
\end{Shaded}

\begin{verbatim}
## 
##    1    2    3    4    5    6 <NA> 
##  149  293  412  826  743  352   25
\end{verbatim}

\begin{Shaded}
\begin{Highlighting}[]
\FunctionTok{cat}\NormalTok{(}\StringTok{"Scale: 1=Very Inaccurate, 6=Very Accurate}\SpecialCharTok{\textbackslash{}n}\StringTok{"}\NormalTok{)}
\end{Highlighting}
\end{Shaded}

\begin{verbatim}
## Scale: 1=Very Inaccurate, 6=Very Accurate
\end{verbatim}

\begin{table}
\centering
\caption{\label{tab:item-table}IPIP Big Five Items. * = reverse-keyed (already recoded in bfi).}
\centering
\begin{tabu} to \linewidth {>{\raggedright}X>{\raggedright}X>{\raggedright}X}
\hline
Item & Trait & Wording\\
\hline
\ttfamily{\textbf{A1}} &  & Am indifferent to the feelings of others.*\\
\cline{1-1}
\cline{3-3}
\ttfamily{\textbf{A2}} &  & Inquire about others' well-being.\\
\cline{1-1}
\cline{3-3}
\ttfamily{\textbf{A3}} &  & Know how to comfort others.\\
\cline{1-1}
\cline{3-3}
\ttfamily{\textbf{A4}} &  & Love children.\\
\cline{1-1}
\cline{3-3}
\ttfamily{\textbf{A5}} & \multirow{-5}{*}{\raggedright\arraybackslash \textcolor[HTML]{2C6FAC}{Agreeableness}} & Make people feel at ease.\\
\cline{1-3}
\ttfamily{\textbf{C1}} &  & Am always prepared.\\
\cline{1-1}
\cline{3-3}
\ttfamily{\textbf{C2}} &  & Leave my belongings around.*\\
\cline{1-1}
\cline{3-3}
\ttfamily{\textbf{C3}} &  & Pay attention to details.\\
\cline{1-1}
\cline{3-3}
\ttfamily{\textbf{C4}} &  & Make a mess of things.*\\
\cline{1-1}
\cline{3-3}
\ttfamily{\textbf{C5}} & \multirow{-5}{*}{\raggedright\arraybackslash \textcolor[HTML]{2C6FAC}{Conscientiousness}} & Often forget to put things back.*\\
\cline{1-3}
\ttfamily{\textbf{E1}} &  & Don't talk a lot.*\\
\cline{1-1}
\cline{3-3}
\ttfamily{\textbf{E2}} &  & Find it difficult to approach others.*\\
\cline{1-1}
\cline{3-3}
\ttfamily{\textbf{E3}} &  & Know how to captivate people.\\
\cline{1-1}
\cline{3-3}
\ttfamily{\textbf{E4}} &  & Make friends easily.\\
\cline{1-1}
\cline{3-3}
\ttfamily{\textbf{E5}} & \multirow{-5}{*}{\raggedright\arraybackslash \textcolor[HTML]{2C6FAC}{Extraversion}} & Take charge.\\
\cline{1-3}
\ttfamily{\textbf{N1}} &  & Get irritated easily.\\
\cline{1-1}
\cline{3-3}
\ttfamily{\textbf{N2}} &  & Seldom feel blue.*\\
\cline{1-1}
\cline{3-3}
\ttfamily{\textbf{N3}} &  & Have frequent mood swings.\\
\cline{1-1}
\cline{3-3}
\ttfamily{\textbf{N4}} &  & Often feel blue.\\
\cline{1-1}
\cline{3-3}
\ttfamily{\textbf{N5}} & \multirow{-5}{*}{\raggedright\arraybackslash \textcolor[HTML]{2C6FAC}{Neuroticism}} & Panic easily.\\
\cline{1-3}
\ttfamily{\textbf{O1}} &  & Am not interested in abstract ideas.*\\
\cline{1-1}
\cline{3-3}
\ttfamily{\textbf{O2}} &  & Have difficulty understanding abstract ideas.*\\
\cline{1-1}
\cline{3-3}
\ttfamily{\textbf{O3}} &  & Do not have a good imagination.*\\
\cline{1-1}
\cline{3-3}
\ttfamily{\textbf{O4}} &  & Have a vivid imagination.\\
\cline{1-1}
\cline{3-3}
\ttfamily{\textbf{O5}} & \multirow{-5}{*}{\raggedright\arraybackslash \textcolor[HTML]{2C6FAC}{Openness}} & Am quick to understand things.\\
\hline
\end{tabu}
\end{table}

\textbf{Reverse-keyed items} are items where agreement indicates
\emph{low} levels of the trait. ``Am indifferent to the feelings of
others'' measures \emph{low} agreeableness. The \texttt{bfi} dataset has
already reverse-coded these items, so higher scores always mean more of
the trait. Always verify this before analysis.

\begin{Shaded}
\begin{Highlighting}[]
\CommentTok{\# Extract the 25 personality items; remove cases with any missing}
\NormalTok{bfi\_items    }\OtherTok{\textless{}{-}}\NormalTok{ bfi[, }\DecValTok{1}\SpecialCharTok{:}\DecValTok{25}\NormalTok{]}
\NormalTok{bfi\_complete }\OtherTok{\textless{}{-}} \FunctionTok{na.omit}\NormalTok{(bfi\_items)}

\FunctionTok{cat}\NormalTok{(}\StringTok{"Total respondents:    "}\NormalTok{, }\FunctionTok{nrow}\NormalTok{(bfi\_items), }\StringTok{"}\SpecialCharTok{\textbackslash{}n}\StringTok{"}\NormalTok{)}
\end{Highlighting}
\end{Shaded}

\begin{verbatim}
## Total respondents:     2800
\end{verbatim}

\begin{Shaded}
\begin{Highlighting}[]
\FunctionTok{cat}\NormalTok{(}\StringTok{"Complete cases:        "}\NormalTok{, }\FunctionTok{nrow}\NormalTok{(bfi\_complete), }\StringTok{"}\SpecialCharTok{\textbackslash{}n}\StringTok{"}\NormalTok{)}
\end{Highlighting}
\end{Shaded}

\begin{verbatim}
## Complete cases:         2436
\end{verbatim}

\begin{Shaded}
\begin{Highlighting}[]
\FunctionTok{cat}\NormalTok{(}\StringTok{"Dropped (any NA):     "}\NormalTok{, }\FunctionTok{nrow}\NormalTok{(bfi\_items) }\SpecialCharTok{{-}} \FunctionTok{nrow}\NormalTok{(bfi\_complete), }\StringTok{"}\SpecialCharTok{\textbackslash{}n}\StringTok{"}\NormalTok{)}
\end{Highlighting}
\end{Shaded}

\begin{verbatim}
## Dropped (any NA):      364
\end{verbatim}

\begin{Shaded}
\begin{Highlighting}[]
\FunctionTok{cat}\NormalTok{(}\StringTok{"Loss rate:            "}\NormalTok{, }\FunctionTok{round}\NormalTok{((}\FunctionTok{nrow}\NormalTok{(bfi\_items)}\SpecialCharTok{{-}}\FunctionTok{nrow}\NormalTok{(bfi\_complete))}\SpecialCharTok{/}\FunctionTok{nrow}\NormalTok{(bfi\_items)}\SpecialCharTok{*}\DecValTok{100}\NormalTok{,}\DecValTok{1}\NormalTok{), }\StringTok{"\%}\SpecialCharTok{\textbackslash{}n}\StringTok{"}\NormalTok{)}
\end{Highlighting}
\end{Shaded}

\begin{verbatim}
## Loss rate:             13 %
\end{verbatim}

\begin{center}\rule{0.5\linewidth}{0.5pt}\end{center}

\section{Part A: Exploratory Factor
Analysis}\label{part-a-exploratory-factor-analysis}

\subsection{Step 1 --- Check
Factorability}\label{step-1-check-factorability}

Before running EFA, check whether the data is suitable. Two standard
tests:

\textbf{KMO (Kaiser-Meyer-Olkin):} Measures whether the correlation
patterns are compact enough to yield reliable factors. Values range from
0 to 1. Interpretation: ≥ .90 = ``Marvelous'', ≥ .80 = ``Meritorious'',
≥ .70 = ``Middling'', ≥ .60 = ``Mediocre'', \textless{} .50 =
``Unacceptable --- do not factor analyse.''

\textbf{Bartlett's test of sphericity:} Tests H₀: the correlation matrix
is an identity matrix (all correlations = 0). A significant result means
correlations are present. With large N this is almost always significant
--- treat it as a minimum sanity check, not a strong criterion.

\begin{Shaded}
\begin{Highlighting}[]
\NormalTok{kmo\_res }\OtherTok{\textless{}{-}} \FunctionTok{KMO}\NormalTok{(bfi\_complete)}

\FunctionTok{cat}\NormalTok{(}\StringTok{"=== KMO Test ===}\SpecialCharTok{\textbackslash{}n}\StringTok{"}\NormalTok{)}
\end{Highlighting}
\end{Shaded}

\begin{verbatim}
## === KMO Test ===
\end{verbatim}

\begin{Shaded}
\begin{Highlighting}[]
\FunctionTok{cat}\NormalTok{(}\StringTok{"Overall MSA:"}\NormalTok{, }\FunctionTok{round}\NormalTok{(kmo\_res}\SpecialCharTok{$}\NormalTok{MSA, }\DecValTok{3}\NormalTok{), }\StringTok{"}\SpecialCharTok{\textbackslash{}n}\StringTok{"}\NormalTok{)}
\end{Highlighting}
\end{Shaded}

\begin{verbatim}
## Overall MSA: 0.849
\end{verbatim}

\begin{Shaded}
\begin{Highlighting}[]
\FunctionTok{cat}\NormalTok{(}\StringTok{"Interpretation: "}\NormalTok{,}
  \FunctionTok{ifelse}\NormalTok{(kmo\_res}\SpecialCharTok{$}\NormalTok{MSA }\SpecialCharTok{\textgreater{}=}\NormalTok{ .}\DecValTok{90}\NormalTok{, }\StringTok{"Marvelous"}\NormalTok{,}
  \FunctionTok{ifelse}\NormalTok{(kmo\_res}\SpecialCharTok{$}\NormalTok{MSA }\SpecialCharTok{\textgreater{}=}\NormalTok{ .}\DecValTok{80}\NormalTok{, }\StringTok{"Meritorious"}\NormalTok{,}
  \FunctionTok{ifelse}\NormalTok{(kmo\_res}\SpecialCharTok{$}\NormalTok{MSA }\SpecialCharTok{\textgreater{}=}\NormalTok{ .}\DecValTok{70}\NormalTok{, }\StringTok{"Middling"}\NormalTok{,}
  \FunctionTok{ifelse}\NormalTok{(kmo\_res}\SpecialCharTok{$}\NormalTok{MSA }\SpecialCharTok{\textgreater{}=}\NormalTok{ .}\DecValTok{60}\NormalTok{, }\StringTok{"Mediocre"}\NormalTok{, }\StringTok{"Unacceptable"}\NormalTok{)))), }\StringTok{"}\SpecialCharTok{\textbackslash{}n\textbackslash{}n}\StringTok{"}\NormalTok{)}
\end{Highlighting}
\end{Shaded}

\begin{verbatim}
## Interpretation:  Meritorious
\end{verbatim}

\begin{Shaded}
\begin{Highlighting}[]
\CommentTok{\# Items with lowest individual KMO (potential weak links)}
\FunctionTok{cat}\NormalTok{(}\StringTok{"Items with lowest individual KMO (flag if \textless{} .50):}\SpecialCharTok{\textbackslash{}n}\StringTok{"}\NormalTok{)}
\end{Highlighting}
\end{Shaded}

\begin{verbatim}
## Items with lowest individual KMO (flag if < .50):
\end{verbatim}

\begin{Shaded}
\begin{Highlighting}[]
\FunctionTok{sort}\NormalTok{(}\FunctionTok{round}\NormalTok{(kmo\_res}\SpecialCharTok{$}\NormalTok{MSAi, }\DecValTok{3}\NormalTok{))[}\DecValTok{1}\SpecialCharTok{:}\DecValTok{5}\NormalTok{] }\SpecialCharTok{|\textgreater{}} \FunctionTok{print}\NormalTok{()}
\end{Highlighting}
\end{Shaded}

\begin{verbatim}
##    A1    O5    O4    N1    N2 
## 0.754 0.762 0.770 0.779 0.780
\end{verbatim}

\begin{Shaded}
\begin{Highlighting}[]
\CommentTok{\# Bartlett}
\NormalTok{bart }\OtherTok{\textless{}{-}} \FunctionTok{cortest.bartlett}\NormalTok{(}\FunctionTok{cor}\NormalTok{(bfi\_complete), }\AttributeTok{n=}\FunctionTok{nrow}\NormalTok{(bfi\_complete))}
\FunctionTok{cat}\NormalTok{(}\StringTok{"}\SpecialCharTok{\textbackslash{}n}\StringTok{=== Bartlett\textquotesingle{}s Test of Sphericity ===}\SpecialCharTok{\textbackslash{}n}\StringTok{"}\NormalTok{)}
\end{Highlighting}
\end{Shaded}

\begin{verbatim}
## 
## === Bartlett's Test of Sphericity ===
\end{verbatim}

\begin{Shaded}
\begin{Highlighting}[]
\FunctionTok{cat}\NormalTok{(}\StringTok{"Chi{-}squared:"}\NormalTok{, }\FunctionTok{round}\NormalTok{(bart}\SpecialCharTok{$}\NormalTok{chisq, }\DecValTok{1}\NormalTok{),}
    \StringTok{"| df:"}\NormalTok{, bart}\SpecialCharTok{$}\NormalTok{df,}
    \StringTok{"| p:"}\NormalTok{, }\FunctionTok{format}\NormalTok{(bart}\SpecialCharTok{$}\NormalTok{p.value, }\AttributeTok{scientific=}\ConstantTok{TRUE}\NormalTok{), }\StringTok{"}\SpecialCharTok{\textbackslash{}n}\StringTok{"}\NormalTok{)}
\end{Highlighting}
\end{Shaded}

\begin{verbatim}
## Chi-squared: 18146.1 | df: 300 | p: 0e+00
\end{verbatim}

\begin{Shaded}
\begin{Highlighting}[]
\FunctionTok{cat}\NormalTok{(}\StringTok{"Conclusion: Correlation matrix is significantly non{-}identity."}\NormalTok{,}
    \StringTok{"Proceed with EFA.}\SpecialCharTok{\textbackslash{}n}\StringTok{"}\NormalTok{)}
\end{Highlighting}
\end{Shaded}

\begin{verbatim}
## Conclusion: Correlation matrix is significantly non-identity. Proceed with EFA.
\end{verbatim}

\subsection{Step 2 --- Determine the Number of
Factors}\label{step-2-determine-the-number-of-factors}

\textbf{Parallel analysis} is the gold standard. It simulates random
datasets of the same size, extracts eigenvalues, and compares: you
retain factors whose eigenvalues exceed what you would expect from
random data.

\begin{Shaded}
\begin{Highlighting}[]
\FunctionTok{set.seed}\NormalTok{(}\DecValTok{2024}\NormalTok{)}
\NormalTok{pa }\OtherTok{\textless{}{-}} \FunctionTok{fa.parallel}\NormalTok{(}
\NormalTok{  bfi\_complete,}
  \AttributeTok{fa      =} \StringTok{"fa"}\NormalTok{,        }\CommentTok{\# factor analysis, not PCA}
  \AttributeTok{fm      =} \StringTok{"minres"}\NormalTok{,    }\CommentTok{\# minimum residual extraction}
  \AttributeTok{n.iter  =} \DecValTok{500}\NormalTok{,         }\CommentTok{\# number of random simulations}
  \AttributeTok{main    =} \StringTok{"Parallel Analysis — IPIP Big Five"}
\NormalTok{)}
\end{Highlighting}
\end{Shaded}

\begin{center}\includegraphics{module_03_efa_cfa_files/figure-latex/parallel-analysis-1} \end{center}

\begin{verbatim}
## Parallel analysis suggests that the number of factors =  6  and the number of components =  NA
\end{verbatim}

\begin{Shaded}
\begin{Highlighting}[]
\FunctionTok{cat}\NormalTok{(}\StringTok{"}\SpecialCharTok{\textbackslash{}n}\StringTok{Parallel analysis recommends:"}\NormalTok{, pa}\SpecialCharTok{$}\NormalTok{nfact, }\StringTok{"factors}\SpecialCharTok{\textbackslash{}n}\StringTok{"}\NormalTok{)}
\end{Highlighting}
\end{Shaded}

\begin{verbatim}
## 
## Parallel analysis recommends: 6 factors
\end{verbatim}

\begin{Shaded}
\begin{Highlighting}[]
\FunctionTok{cat}\NormalTok{(}\StringTok{"}\SpecialCharTok{\textbackslash{}n}\StringTok{How to read the plot:}\SpecialCharTok{\textbackslash{}n}\StringTok{"}\NormalTok{)}
\end{Highlighting}
\end{Shaded}

\begin{verbatim}
## 
## How to read the plot:
\end{verbatim}

\begin{Shaded}
\begin{Highlighting}[]
\FunctionTok{cat}\NormalTok{(}\StringTok{"  Solid line = your actual eigenvalues}\SpecialCharTok{\textbackslash{}n}\StringTok{"}\NormalTok{)}
\end{Highlighting}
\end{Shaded}

\begin{verbatim}
##   Solid line = your actual eigenvalues
\end{verbatim}

\begin{Shaded}
\begin{Highlighting}[]
\FunctionTok{cat}\NormalTok{(}\StringTok{"  Dashed line = eigenvalues from random (simulated) data}\SpecialCharTok{\textbackslash{}n}\StringTok{"}\NormalTok{)}
\end{Highlighting}
\end{Shaded}

\begin{verbatim}
##   Dashed line = eigenvalues from random (simulated) data
\end{verbatim}

\begin{Shaded}
\begin{Highlighting}[]
\FunctionTok{cat}\NormalTok{(}\StringTok{"  Retain factors where solid \textgreater{} dashed (to the left of where lines cross)}\SpecialCharTok{\textbackslash{}n}\StringTok{"}\NormalTok{)}
\end{Highlighting}
\end{Shaded}

\begin{verbatim}
##   Retain factors where solid > dashed (to the left of where lines cross)
\end{verbatim}

\textbf{Why not the eigenvalue \textgreater{} 1 rule?} The classic
Kaiser criterion says retain factors with eigenvalue \textgreater{} 1.
With 25 items this typically suggests 6--8 factors --- almost always too
many. The rule has no theoretical basis and is known to over-extract.
Use parallel analysis instead.

\subsection{Step 3 --- Run EFA}\label{step-3-run-efa}

Key decisions: - \textbf{Extraction method:} We use \texttt{minres}
(minimum residual) --- it minimises off-diagonal residuals, does not
assume multivariate normality, and is robust to non-normality. -
\textbf{Rotation:} We use \texttt{oblimin} (oblique) --- it allows
factors to correlate, which is realistic for personality traits and most
social science constructs. If you assume traits are uncorrelated, use
\texttt{varimax} (orthogonal).

\begin{Shaded}
\begin{Highlighting}[]
\CommentTok{\# Fit 5{-}factor EFA (expected Big Five solution)}
\NormalTok{efa5 }\OtherTok{\textless{}{-}} \FunctionTok{fa}\NormalTok{(}
  \AttributeTok{r         =}\NormalTok{ bfi\_complete,}
  \AttributeTok{nfactors  =} \DecValTok{5}\NormalTok{,}
  \AttributeTok{rotate    =} \StringTok{"oblimin"}\NormalTok{,    }\CommentTok{\# oblique: factors can correlate}
  \AttributeTok{fm        =} \StringTok{"minres"}\NormalTok{,}
  \AttributeTok{scores    =} \StringTok{"regression"}  \CommentTok{\# compute factor scores for each respondent}
\NormalTok{)}

\CommentTok{\# Print full output with cut = 0.30 (suppress loadings below |.30| for clarity)}
\FunctionTok{print}\NormalTok{(efa5, }\AttributeTok{cut =} \FloatTok{0.30}\NormalTok{, }\AttributeTok{sort =} \ConstantTok{TRUE}\NormalTok{, }\AttributeTok{digits =} \DecValTok{2}\NormalTok{)}
\end{Highlighting}
\end{Shaded}

\begin{verbatim}
## Factor Analysis using method =  minres
## Call: fa(r = bfi_complete, nfactors = 5, rotate = "oblimin", scores = "regression", 
##     fm = "minres")
## Standardized loadings (pattern matrix) based upon correlation matrix
##    item   MR2   MR1   MR3   MR5   MR4   h2   u2 com
## N1   16  0.83                         0.68 0.32 1.1
## N2   17  0.78                         0.61 0.39 1.0
## N3   18  0.70                         0.54 0.46 1.1
## N5   20  0.48                         0.35 0.65 2.0
## N4   19  0.47  0.40                   0.51 0.49 2.3
## E2   12        0.67                   0.55 0.45 1.1
## E4   14       -0.59                   0.54 0.46 1.5
## E1   11        0.55                   0.35 0.65 1.2
## E5   15       -0.42                   0.41 0.59 2.7
## E3   13       -0.41              0.30 0.44 0.56 2.7
## C2    7              0.67             0.45 0.55 1.2
## C4    9             -0.64             0.48 0.52 1.1
## C3    8              0.57             0.32 0.68 1.1
## C5   10             -0.56             0.44 0.56 1.4
## C1    6              0.55             0.35 0.65 1.2
## A3    3                    0.68       0.54 0.46 1.1
## A2    2                    0.66       0.46 0.54 1.0
## A5    5                    0.54       0.47 0.53 1.5
## A4    4                    0.45       0.30 0.70 1.7
## A1    1                   -0.44       0.20 0.80 1.9
## O3   23                          0.62 0.47 0.53 1.2
## O5   25                         -0.54 0.30 0.70 1.2
## O1   21                          0.52 0.32 0.68 1.1
## O2   22                         -0.47 0.27 0.73 1.7
## O4   24        0.33              0.36 0.25 0.75 2.6
## 
##                        MR2  MR1  MR3  MR5  MR4
## SS loadings           2.58 2.21 2.10 2.07 1.64
## Proportion Var        0.10 0.09 0.08 0.08 0.07
## Cumulative Var        0.10 0.19 0.28 0.36 0.42
## Proportion Explained  0.24 0.21 0.20 0.20 0.15
## Cumulative Proportion 0.24 0.45 0.65 0.85 1.00
## 
##  With factor correlations of 
##       MR2   MR1   MR3   MR5   MR4
## MR2  1.00  0.22 -0.19 -0.05  0.00
## MR1  0.22  1.00 -0.24 -0.33 -0.17
## MR3 -0.19 -0.24  1.00  0.20  0.20
## MR5 -0.05 -0.33  0.20  1.00  0.20
## MR4  0.00 -0.17  0.20  0.20  1.00
## 
## Mean item complexity =  1.5
## Test of the hypothesis that 5 factors are sufficient.
## 
## df null model =  300  with the objective function =  7.48 with Chi Square =  18146.07
## df of  the model are 185  and the objective function was  0.64 
## 
## The root mean square of the residuals (RMSR) is  0.03 
## The df corrected root mean square of the residuals is  0.04 
## 
## The harmonic n.obs is  2436 with the empirical chi square  565.95  with prob <  3.1e-40 
## The total n.obs was  2436  with Likelihood Chi Square =  1538.52  with prob <  8.6e-212 
## 
## Tucker Lewis Index of factoring reliability =  0.877
## RMSEA index =  0.055  and the 90 % confidence intervals are  0.052 0.057
## BIC =  95.87
## Fit based upon off diagonal values = 0.98
## Measures of factor score adequacy             
##                                                    MR2  MR1  MR3  MR5  MR4
## Correlation of (regression) scores with factors   0.91 0.82 0.84 0.83 0.82
## Multiple R square of scores with factors          0.83 0.67 0.71 0.68 0.66
## Minimum correlation of possible factor scores     0.66 0.34 0.43 0.36 0.33
\end{verbatim}

\subsection{Reading the Full EFA Output --- Line by
Line}\label{reading-the-full-efa-output-line-by-line}

The \texttt{psych::fa()} output has several sections. Here is what each
means:

\begin{table}
\centering
\caption{\label{tab:efa-output-guide}Complete Guide to psych::fa() Output Sections}
\centering
\begin{tabu} to \linewidth {>{\raggedright}X>{\raggedright}X}
\hline
Output Section & What It Means\\
\hline
\ttfamily{\textbf{Standardised Loadings (Pattern Matrix)}} & The core output. Each value = how strongly that item loads on that factor after rotation. Focus on values > |.30|.\\
\hline
\ttfamily{\textbf{h2 column}} & Communality: proportion of item variance explained by ALL factors combined. If h2 < .30 the item is poorly explained.\\
\hline
\ttfamily{\textbf{u2 column}} & Uniqueness = 1 - h2. The proportion of variance in the item NOT explained by any factor.\\
\hline
\ttfamily{\textbf{com column}} & Complexity: number of factors the item loads substantially on. Ideal = 1.0 (simple structure). Values > 2 signal cross-loaders.\\
\hline
\ttfamily{\textbf{SS loadings row}} & Sum of squared loadings per factor. Larger = more important factor. Computed AFTER rotation.\\
\hline
\ttfamily{\textbf{Proportion Var row}} & SS loadings divided by total items. Proportion of total variance this factor explains.\\
\hline
\ttfamily{\textbf{Cumulative Var row}} & Running total of Proportion Var across factors. At 5 factors: total explained variance.\\
\hline
\ttfamily{\textbf{Proportion Explained row}} & Each factor's proportion of the variance that IS explained (not of total variance).\\
\hline
\ttfamily{\textbf{Mean item complexity}} & Average complexity across items. Values near 1 indicate clean simple structure.\\
\hline
\ttfamily{\textbf{RMSEA / TLI}} & These are model fit indices for the EFA as a structural model. Interpreted same as CFA.\\
\hline
\ttfamily{\textbf{Root mean square of residuals}} & Average absolute residual correlation. Should be < .05 for good fit.\\
\hline
\ttfamily{\textbf{Fit based upon off diagonal}} & χ² test on off-diagonal residuals. Significant = model does not perfectly reproduce correlations.\\
\hline
\end{tabu}
\end{table}

\subsection{Interpreting the Loadings}\label{interpreting-the-loadings}

\begin{Shaded}
\begin{Highlighting}[]
\CommentTok{\# Extract loadings cleanly for annotation}
\NormalTok{L }\OtherTok{\textless{}{-}} \FunctionTok{unclass}\NormalTok{(efa5}\SpecialCharTok{$}\NormalTok{loadings)}

\CommentTok{\# Show the Agreeableness items as a worked example}
\FunctionTok{cat}\NormalTok{(}\StringTok{"=== Worked example: Agreeableness items (A1–A5) ===}\SpecialCharTok{\textbackslash{}n\textbackslash{}n}\StringTok{"}\NormalTok{)}
\end{Highlighting}
\end{Shaded}

\begin{verbatim}
## === Worked example: Agreeableness items (A1–A5) ===
\end{verbatim}

\begin{Shaded}
\begin{Highlighting}[]
\FunctionTok{cat}\NormalTok{(}\StringTok{"Factor loading table for A items:}\SpecialCharTok{\textbackslash{}n}\StringTok{"}\NormalTok{)}
\end{Highlighting}
\end{Shaded}

\begin{verbatim}
## Factor loading table for A items:
\end{verbatim}

\begin{Shaded}
\begin{Highlighting}[]
\FunctionTok{round}\NormalTok{(L[}\FunctionTok{paste0}\NormalTok{(}\StringTok{"A"}\NormalTok{,}\DecValTok{1}\SpecialCharTok{:}\DecValTok{5}\NormalTok{), ], }\DecValTok{2}\NormalTok{) }\SpecialCharTok{|\textgreater{}} \FunctionTok{print}\NormalTok{()}
\end{Highlighting}
\end{Shaded}

\begin{verbatim}
##      MR2   MR1  MR3   MR5   MR4
## A1  0.20 -0.18 0.07 -0.44 -0.05
## A2 -0.02  0.00 0.07  0.66  0.03
## A3 -0.02 -0.10 0.03  0.68  0.03
## A4 -0.05 -0.08 0.19  0.45 -0.16
## A5 -0.12 -0.23 0.00  0.54  0.05
\end{verbatim}

\begin{Shaded}
\begin{Highlighting}[]
\FunctionTok{cat}\NormalTok{(}\StringTok{"}\SpecialCharTok{\textbackslash{}n}\StringTok{h2 (communality) for A items:}\SpecialCharTok{\textbackslash{}n}\StringTok{"}\NormalTok{)}
\end{Highlighting}
\end{Shaded}

\begin{verbatim}
## 
## h2 (communality) for A items:
\end{verbatim}

\begin{Shaded}
\begin{Highlighting}[]
\FunctionTok{round}\NormalTok{(efa5}\SpecialCharTok{$}\NormalTok{communality[}\FunctionTok{paste0}\NormalTok{(}\StringTok{"A"}\NormalTok{,}\DecValTok{1}\SpecialCharTok{:}\DecValTok{5}\NormalTok{)], }\DecValTok{2}\NormalTok{) }\SpecialCharTok{|\textgreater{}} \FunctionTok{print}\NormalTok{()}
\end{Highlighting}
\end{Shaded}

\begin{verbatim}
##   A1   A2   A3   A4   A5 
## 0.20 0.46 0.54 0.30 0.47
\end{verbatim}

\begin{Shaded}
\begin{Highlighting}[]
\FunctionTok{cat}\NormalTok{(}\StringTok{"}\SpecialCharTok{\textbackslash{}n}\StringTok{u2 (uniqueness) for A items:}\SpecialCharTok{\textbackslash{}n}\StringTok{"}\NormalTok{)}
\end{Highlighting}
\end{Shaded}

\begin{verbatim}
## 
## u2 (uniqueness) for A items:
\end{verbatim}

\begin{Shaded}
\begin{Highlighting}[]
\FunctionTok{round}\NormalTok{(efa5}\SpecialCharTok{$}\NormalTok{uniquenesses[}\FunctionTok{paste0}\NormalTok{(}\StringTok{"A"}\NormalTok{,}\DecValTok{1}\SpecialCharTok{:}\DecValTok{5}\NormalTok{)], }\DecValTok{2}\NormalTok{) }\SpecialCharTok{|\textgreater{}} \FunctionTok{print}\NormalTok{()}
\end{Highlighting}
\end{Shaded}

\begin{verbatim}
##   A1   A2   A3   A4   A5 
## 0.80 0.54 0.46 0.70 0.53
\end{verbatim}

\begin{Shaded}
\begin{Highlighting}[]
\FunctionTok{cat}\NormalTok{(}\StringTok{"}\SpecialCharTok{\textbackslash{}n}\StringTok{What this means:}\SpecialCharTok{\textbackslash{}n}\StringTok{"}\NormalTok{)}
\end{Highlighting}
\end{Shaded}

\begin{verbatim}
## 
## What this means:
\end{verbatim}

\begin{Shaded}
\begin{Highlighting}[]
\FunctionTok{cat}\NormalTok{(}\StringTok{"  A1 loads .XX on its factor — this is the \textquotesingle{}primary loading\textquotesingle{}.}\SpecialCharTok{\textbackslash{}n}\StringTok{"}\NormalTok{)}
\end{Highlighting}
\end{Shaded}

\begin{verbatim}
##   A1 loads .XX on its factor — this is the 'primary loading'.
\end{verbatim}

\begin{Shaded}
\begin{Highlighting}[]
\FunctionTok{cat}\NormalTok{(}\StringTok{"  h2 = .XX means the factors explain XX\% of A1\textquotesingle{}s total variance.}\SpecialCharTok{\textbackslash{}n}\StringTok{"}\NormalTok{)}
\end{Highlighting}
\end{Shaded}

\begin{verbatim}
##   h2 = .XX means the factors explain XX% of A1's total variance.
\end{verbatim}

\begin{Shaded}
\begin{Highlighting}[]
\FunctionTok{cat}\NormalTok{(}\StringTok{"  u2 = .XX means XX\% of A1\textquotesingle{}s variance is item{-}specific (measurement error + unique content).}\SpecialCharTok{\textbackslash{}n}\StringTok{"}\NormalTok{)}
\end{Highlighting}
\end{Shaded}

\begin{verbatim}
##   u2 = .XX means XX% of A1's variance is item-specific (measurement error + unique content).
\end{verbatim}

\begin{Shaded}
\begin{Highlighting}[]
\FunctionTok{cat}\NormalTok{(}\StringTok{"  If A1 has h2 \textless{} .30, consider removing it from the scale.}\SpecialCharTok{\textbackslash{}n}\StringTok{"}\NormalTok{)}
\end{Highlighting}
\end{Shaded}

\begin{verbatim}
##   If A1 has h2 < .30, consider removing it from the scale.
\end{verbatim}

\begin{Shaded}
\begin{Highlighting}[]
\CommentTok{\# In oblique rotation, factors are allowed to correlate}
\CommentTok{\# The factor correlation matrix tells you HOW correlated they are}
\FunctionTok{cat}\NormalTok{(}\StringTok{"=== Factor Correlation Matrix (Phi) ===}\SpecialCharTok{\textbackslash{}n}\StringTok{"}\NormalTok{)}
\end{Highlighting}
\end{Shaded}

\begin{verbatim}
## === Factor Correlation Matrix (Phi) ===
\end{verbatim}

\begin{Shaded}
\begin{Highlighting}[]
\FunctionTok{cat}\NormalTok{(}\StringTok{"(Only available with oblique rotation — orthogonal rotation forces these to 0)}\SpecialCharTok{\textbackslash{}n\textbackslash{}n}\StringTok{"}\NormalTok{)}
\end{Highlighting}
\end{Shaded}

\begin{verbatim}
## (Only available with oblique rotation — orthogonal rotation forces these to 0)
\end{verbatim}

\begin{Shaded}
\begin{Highlighting}[]
\FunctionTok{round}\NormalTok{(efa5}\SpecialCharTok{$}\NormalTok{Phi, }\DecValTok{3}\NormalTok{) }\SpecialCharTok{|\textgreater{}} \FunctionTok{print}\NormalTok{()}
\end{Highlighting}
\end{Shaded}

\begin{verbatim}
##        MR2    MR1    MR3    MR5    MR4
## MR2  1.000  0.217 -0.191 -0.046 -0.001
## MR1  0.217  1.000 -0.237 -0.330 -0.166
## MR3 -0.191 -0.237  1.000  0.202  0.198
## MR5 -0.046 -0.330  0.202  1.000  0.196
## MR4 -0.001 -0.166  0.198  0.196  1.000
\end{verbatim}

\begin{Shaded}
\begin{Highlighting}[]
\FunctionTok{cat}\NormalTok{(}\StringTok{"}\SpecialCharTok{\textbackslash{}n}\StringTok{Interpretation:}\SpecialCharTok{\textbackslash{}n}\StringTok{"}\NormalTok{)}
\end{Highlighting}
\end{Shaded}

\begin{verbatim}
## 
## Interpretation:
\end{verbatim}

\begin{Shaded}
\begin{Highlighting}[]
\FunctionTok{cat}\NormalTok{(}\StringTok{"  Values near 0: factors are nearly orthogonal (empirically independent)}\SpecialCharTok{\textbackslash{}n}\StringTok{"}\NormalTok{)}
\end{Highlighting}
\end{Shaded}

\begin{verbatim}
##   Values near 0: factors are nearly orthogonal (empirically independent)
\end{verbatim}

\begin{Shaded}
\begin{Highlighting}[]
\FunctionTok{cat}\NormalTok{(}\StringTok{"  Values \textgreater{} |.30|: factors are meaningfully correlated (use oblique, not orthogonal)}\SpecialCharTok{\textbackslash{}n}\StringTok{"}\NormalTok{)}
\end{Highlighting}
\end{Shaded}

\begin{verbatim}
##   Values > |.30|: factors are meaningfully correlated (use oblique, not orthogonal)
\end{verbatim}

\begin{Shaded}
\begin{Highlighting}[]
\FunctionTok{cat}\NormalTok{(}\StringTok{"  Values \textgreater{} |.70|: factors may be too similar — consider merging them}\SpecialCharTok{\textbackslash{}n}\StringTok{"}\NormalTok{)}
\end{Highlighting}
\end{Shaded}

\begin{verbatim}
##   Values > |.70|: factors may be too similar — consider merging them
\end{verbatim}

\subsection{Visualising the Factor
Structure}\label{visualising-the-factor-structure}

\begin{Shaded}
\begin{Highlighting}[]
\NormalTok{L\_df }\OtherTok{\textless{}{-}} \FunctionTok{as.data.frame}\NormalTok{(L) }\SpecialCharTok{|\textgreater{}}
  \FunctionTok{rownames\_to\_column}\NormalTok{(}\StringTok{"Item"}\NormalTok{) }\SpecialCharTok{|\textgreater{}}
  \FunctionTok{mutate}\NormalTok{(}
    \AttributeTok{Trait =} \FunctionTok{str\_extract}\NormalTok{(Item, }\StringTok{"\^{}[A{-}Z]"}\NormalTok{),}
    \AttributeTok{Trait =}\NormalTok{ dplyr}\SpecialCharTok{::}\FunctionTok{recode}\NormalTok{(Trait, }\AttributeTok{A=}\StringTok{"Agree."}\NormalTok{, }\AttributeTok{C=}\StringTok{"Consc."}\NormalTok{, }\AttributeTok{E=}\StringTok{"Extra."}\NormalTok{, }\AttributeTok{N=}\StringTok{"Neuro."}\NormalTok{, }\AttributeTok{O=}\StringTok{"Open."}\NormalTok{)}
\NormalTok{  ) }\SpecialCharTok{|\textgreater{}}
  \FunctionTok{pivot\_longer}\NormalTok{(}\FunctionTok{starts\_with}\NormalTok{(}\StringTok{"MR"}\NormalTok{), }\AttributeTok{names\_to=}\StringTok{"Factor"}\NormalTok{, }\AttributeTok{values\_to=}\StringTok{"Loading"}\NormalTok{)}

\CommentTok{\# Map factor names based on which items load highest}
\NormalTok{factor\_labels }\OtherTok{\textless{}{-}}\NormalTok{ L\_df }\SpecialCharTok{|\textgreater{}}
  \FunctionTok{group\_by}\NormalTok{(Factor) }\SpecialCharTok{|\textgreater{}}
  \FunctionTok{slice\_max}\NormalTok{(}\FunctionTok{abs}\NormalTok{(Loading), }\AttributeTok{n=}\DecValTok{3}\NormalTok{) }\SpecialCharTok{|\textgreater{}}
  \FunctionTok{summarise}\NormalTok{(}\AttributeTok{top\_trait =} \FunctionTok{names}\NormalTok{(}\FunctionTok{which.max}\NormalTok{(}\FunctionTok{table}\NormalTok{(}\FunctionTok{str\_extract}\NormalTok{(Item,}\StringTok{"\^{}[A{-}Z]"}\NormalTok{)))),}
            \AttributeTok{.groups=}\StringTok{"drop"}\NormalTok{) }\SpecialCharTok{|\textgreater{}}
  \FunctionTok{mutate}\NormalTok{(}\AttributeTok{Label =}\NormalTok{ dplyr}\SpecialCharTok{::}\FunctionTok{recode}\NormalTok{(top\_trait,}
    \AttributeTok{A=}\StringTok{"F: Agreeableness"}\NormalTok{, }\AttributeTok{C=}\StringTok{"F: Conscientiousness"}\NormalTok{,}
    \AttributeTok{E=}\StringTok{"F: Extraversion"}\NormalTok{,  }\AttributeTok{N=}\StringTok{"F: Neuroticism"}\NormalTok{, }\AttributeTok{O=}\StringTok{"F: Openness"}\NormalTok{))}

\NormalTok{L\_df }\OtherTok{\textless{}{-}}\NormalTok{ L\_df }\SpecialCharTok{|\textgreater{}}
  \FunctionTok{left\_join}\NormalTok{(factor\_labels, }\AttributeTok{by=}\StringTok{"Factor"}\NormalTok{)}

\FunctionTok{ggplot}\NormalTok{(L\_df, }\FunctionTok{aes}\NormalTok{(}\AttributeTok{x=}\NormalTok{Label, }\AttributeTok{y=}\FunctionTok{fct\_rev}\NormalTok{(Item), }\AttributeTok{fill=}\NormalTok{Loading)) }\SpecialCharTok{+}
  \FunctionTok{geom\_tile}\NormalTok{(}\AttributeTok{colour=}\StringTok{"white"}\NormalTok{, }\AttributeTok{linewidth=}\FloatTok{0.4}\NormalTok{) }\SpecialCharTok{+}
  \FunctionTok{geom\_text}\NormalTok{(}\FunctionTok{aes}\NormalTok{(}\AttributeTok{label=}\FunctionTok{ifelse}\NormalTok{(}\FunctionTok{abs}\NormalTok{(Loading) }\SpecialCharTok{\textgreater{}=} \FloatTok{0.25}\NormalTok{, }\FunctionTok{round}\NormalTok{(Loading,}\DecValTok{2}\NormalTok{), }\StringTok{""}\NormalTok{),}
                \AttributeTok{colour=}\FunctionTok{abs}\NormalTok{(Loading) }\SpecialCharTok{\textgreater{}} \FloatTok{0.4}\NormalTok{),}
            \AttributeTok{size=}\FloatTok{2.9}\NormalTok{) }\SpecialCharTok{+}
  \FunctionTok{scale\_fill\_gradient2}\NormalTok{(}\AttributeTok{low=}\StringTok{"\#C0392B"}\NormalTok{, }\AttributeTok{mid=}\StringTok{"white"}\NormalTok{, }\AttributeTok{high=}\StringTok{"\#2C6FAC"}\NormalTok{,}
                       \AttributeTok{midpoint=}\DecValTok{0}\NormalTok{, }\AttributeTok{limits=}\FunctionTok{c}\NormalTok{(}\SpecialCharTok{{-}}\DecValTok{1}\NormalTok{,}\DecValTok{1}\NormalTok{), }\AttributeTok{name=}\StringTok{"Loading"}\NormalTok{) }\SpecialCharTok{+}
  \FunctionTok{scale\_colour\_manual}\NormalTok{(}\AttributeTok{values=}\FunctionTok{c}\NormalTok{(}\StringTok{"FALSE"}\OtherTok{=}\StringTok{"grey40"}\NormalTok{,}\StringTok{"TRUE"}\OtherTok{=}\StringTok{"white"}\NormalTok{), }\AttributeTok{guide=}\StringTok{"none"}\NormalTok{) }\SpecialCharTok{+}
  \FunctionTok{facet\_grid}\NormalTok{(Trait}\SpecialCharTok{\textasciitilde{}}\NormalTok{., }\AttributeTok{scales=}\StringTok{"free\_y"}\NormalTok{, }\AttributeTok{space=}\StringTok{"free\_y"}\NormalTok{, }\AttributeTok{switch=}\StringTok{"y"}\NormalTok{) }\SpecialCharTok{+}
  \FunctionTok{labs}\NormalTok{(}
    \AttributeTok{title    =} \StringTok{"EFA Factor Loading Matrix (Oblimin rotation, 5 factors)"}\NormalTok{,}
    \AttributeTok{subtitle =} \FunctionTok{paste}\NormalTok{(}\StringTok{"Loadings ≥ |.25| shown. Strong primary loadings appear as dark blue/red blocks."}\NormalTok{,}
                     \StringTok{"}\SpecialCharTok{\textbackslash{}n}\StringTok{Off{-}diagonal colours = cross{-}loadings (problematic if \textgreater{} |.30|)."}\NormalTok{),}
    \AttributeTok{x=}\StringTok{"Factor"}\NormalTok{, }\AttributeTok{y=}\ConstantTok{NULL}\NormalTok{,}
    \AttributeTok{caption  =} \FunctionTok{paste}\NormalTok{(}\StringTok{"Source: psych::bfi dataset. N ="}\NormalTok{, }\FunctionTok{nrow}\NormalTok{(bfi\_complete), }\StringTok{"."}\NormalTok{)}
\NormalTok{  ) }\SpecialCharTok{+}
  \FunctionTok{theme}\NormalTok{(}\AttributeTok{axis.text.x=}\FunctionTok{element\_text}\NormalTok{(}\AttributeTok{angle=}\DecValTok{30}\NormalTok{, }\AttributeTok{hjust=}\DecValTok{1}\NormalTok{),}
        \AttributeTok{strip.text.y.left=}\FunctionTok{element\_text}\NormalTok{(}\AttributeTok{angle=}\DecValTok{0}\NormalTok{, }\AttributeTok{hjust=}\DecValTok{1}\NormalTok{, }\AttributeTok{face=}\StringTok{"bold"}\NormalTok{),}
        \AttributeTok{strip.placement=}\StringTok{"outside"}\NormalTok{)}
\end{Highlighting}
\end{Shaded}

\begin{center}\includegraphics{module_03_efa_cfa_files/figure-latex/efa-heatmap-1} \end{center}

\textbf{Reading the heatmap:} Each row is an item; each column is a
factor. Dark blue in a cell = strong positive loading. Dark red = strong
negative loading. White = near-zero loading. A clean five-factor
solution shows five ``hot spots'' on the diagonal --- each block of
items (A1--A5, C1--C5, etc.) loads strongly on one factor and weakly on
all others. Cross-loadings (substantial colours off the diagonal)
indicate items that do not fit cleanly into the proposed structure.

\begin{center}\rule{0.5\linewidth}{0.5pt}\end{center}

\section{Part B: Confirmatory Factor
Analysis}\label{part-b-confirmatory-factor-analysis}

\subsection{Moving from EFA to CFA}\label{moving-from-efa-to-cfa}

EFA discovers factor structure from data. CFA tests a
\emph{theoretically specified} factor structure. In CFA you explicitly
declare:

\begin{itemize}
\tightlist
\item
  Which items load on which factors (all other loadings are fixed to
  zero)
\item
  Whether factors are allowed to correlate (usually yes --- lavaan does
  this by default)
\item
  Which residuals are allowed to correlate (usually none, unless
  theoretically justified)
\end{itemize}

This is far more constrained than EFA, so CFA will almost always fit the
data less well than EFA. That is the point: CFA tests whether a specific
theory-driven structure is adequate, not whether it is the best possible
fitting structure.

\textbf{Important rule:} Use EFA on one half of your data; CFA on the
other half. Using EFA to discover a structure and CFA to ``confirm'' it
on the \emph{same data} is circular --- you are confirming what the data
already told you.

\subsection{Specifying the CFA Model}\label{specifying-the-cfa-model}

lavaan syntax uses \texttt{=\textasciitilde{}} (is measured by) to
define latent variables:

\begin{Shaded}
\begin{Highlighting}[]
\NormalTok{big5\_model }\OtherTok{\textless{}{-}} \StringTok{\textquotesingle{}}
\StringTok{  \# MEASUREMENT MODEL}
\StringTok{  \# Each line: LatentVariable =\textasciitilde{} item1 + item2 + item3 ...}
\StringTok{  \# lavaan will estimate one loading per item.}
\StringTok{  \# The first item in each line gets its loading fixed to 1.0 for scaling}
\StringTok{  \# (unless you set std.lv = TRUE, which fixes the factor variance to 1 instead)}

\StringTok{  Agreeableness     =\textasciitilde{} A1 + A2 + A3 + A4 + A5}
\StringTok{  Conscientiousness =\textasciitilde{} C1 + C2 + C3 + C4 + C5}
\StringTok{  Extraversion      =\textasciitilde{} E1 + E2 + E3 + E4 + E5}
\StringTok{  Neuroticism       =\textasciitilde{} N1 + N2 + N3 + N4 + N5}
\StringTok{  Openness          =\textasciitilde{} O1 + O2 + O3 + O4 + O5}

\StringTok{  \# By default lavaan estimates ALL pairwise factor covariances.}
\StringTok{  \# You will see these in the output as: Agreeableness \textasciitilde{}\textasciitilde{} Conscientiousness etc.}
\StringTok{  \# To constrain two specific factors to be uncorrelated, add:}
\StringTok{  \# Agreeableness \textasciitilde{}\textasciitilde{} 0*Conscientiousness}
\StringTok{\textquotesingle{}}
\FunctionTok{cat}\NormalTok{(}\StringTok{"Model specified. Contains:}\SpecialCharTok{\textbackslash{}n}\StringTok{"}\NormalTok{)}
\end{Highlighting}
\end{Shaded}

\begin{verbatim}
## Model specified. Contains:
\end{verbatim}

\begin{Shaded}
\begin{Highlighting}[]
\FunctionTok{cat}\NormalTok{(}\StringTok{"  5 latent variables}\SpecialCharTok{\textbackslash{}n}\StringTok{"}\NormalTok{)}
\end{Highlighting}
\end{Shaded}

\begin{verbatim}
##   5 latent variables
\end{verbatim}

\begin{Shaded}
\begin{Highlighting}[]
\FunctionTok{cat}\NormalTok{(}\StringTok{"  25 observed items (5 per factor)}\SpecialCharTok{\textbackslash{}n}\StringTok{"}\NormalTok{)}
\end{Highlighting}
\end{Shaded}

\begin{verbatim}
##   25 observed items (5 per factor)
\end{verbatim}

\begin{Shaded}
\begin{Highlighting}[]
\FunctionTok{cat}\NormalTok{(}\StringTok{"  5×4/2 = 10 estimated factor covariances}\SpecialCharTok{\textbackslash{}n}\StringTok{"}\NormalTok{)}
\end{Highlighting}
\end{Shaded}

\begin{verbatim}
##   5×4/2 = 10 estimated factor covariances
\end{verbatim}

\begin{Shaded}
\begin{Highlighting}[]
\FunctionTok{cat}\NormalTok{(}\StringTok{"  25 residual variances}\SpecialCharTok{\textbackslash{}n}\StringTok{"}\NormalTok{)}
\end{Highlighting}
\end{Shaded}

\begin{verbatim}
##   25 residual variances
\end{verbatim}

\begin{Shaded}
\begin{Highlighting}[]
\FunctionTok{cat}\NormalTok{(}\StringTok{"  20 estimated loadings (5 factors × 5 items − 5 fixed for scaling)}\SpecialCharTok{\textbackslash{}n}\StringTok{"}\NormalTok{)}
\end{Highlighting}
\end{Shaded}

\begin{verbatim}
##   20 estimated loadings (5 factors × 5 items − 5 fixed for scaling)
\end{verbatim}

\subsection{Fitting the CFA}\label{fitting-the-cfa}

\begin{Shaded}
\begin{Highlighting}[]
\NormalTok{cfa\_fit }\OtherTok{\textless{}{-}} \FunctionTok{cfa}\NormalTok{(}
  \AttributeTok{model     =}\NormalTok{ big5\_model,}
  \AttributeTok{data      =}\NormalTok{ bfi\_complete,}
  \AttributeTok{estimator =} \StringTok{"MLR"}\NormalTok{,    }\CommentTok{\# Maximum Likelihood with Robust standard errors}
                        \CommentTok{\# MLR is preferred over ML when data is non{-}normal}
  \AttributeTok{std.lv    =} \ConstantTok{TRUE}      \CommentTok{\# Fix latent variance to 1 so ALL loadings are estimated}
                        \CommentTok{\# (not just 4 of the 5 per factor)}
\NormalTok{)}

\CommentTok{\# The full summary — always start here}
\FunctionTok{summary}\NormalTok{(cfa\_fit, }\AttributeTok{fit.measures=}\ConstantTok{TRUE}\NormalTok{, }\AttributeTok{standardized=}\ConstantTok{TRUE}\NormalTok{, }\AttributeTok{ci=}\ConstantTok{TRUE}\NormalTok{)}
\end{Highlighting}
\end{Shaded}

\begin{verbatim}
## lavaan 0.6-21 ended normally after 23 iterations
## 
##   Estimator                                         ML
##   Optimization method                           NLMINB
##   Number of model parameters                        60
## 
##   Number of observations                          2436
## 
## Model Test User Model:
##                                               Standard      Scaled
##   Test Statistic                              4165.467    3612.178
##   Degrees of freedom                               265         265
##   P-value (Chi-square)                           0.000       0.000
##   Scaling correction factor                                  1.153
##     Yuan-Bentler correction (Mplus variant)                       
## 
## Model Test Baseline Model:
## 
##   Test statistic                             18222.116   15289.285
##   Degrees of freedom                               300         300
##   P-value                                        0.000       0.000
##   Scaling correction factor                                  1.192
## 
## User Model versus Baseline Model:
## 
##   Comparative Fit Index (CFI)                    0.782       0.777
##   Tucker-Lewis Index (TLI)                       0.754       0.747
##                                                                   
##   Robust Comparative Fit Index (CFI)                         0.784
##   Robust Tucker-Lewis Index (TLI)                            0.755
## 
## Loglikelihood and Information Criteria:
## 
##   Loglikelihood user model (H0)             -99840.238  -99840.238
##   Scaling correction factor                                  1.216
##       for the MLR correction                                      
##   Loglikelihood unrestricted model (H1)     -97757.504  -97757.504
##   Scaling correction factor                                  1.165
##       for the MLR correction                                      
##                                                                   
##   Akaike (AIC)                              199800.476  199800.476
##   Bayesian (BIC)                            200148.363  200148.363
##   Sample-size adjusted Bayesian (SABIC)     199957.729  199957.729
## 
## Root Mean Square Error of Approximation:
## 
##   RMSEA                                          0.078       0.072
##   90 Percent confidence interval - lower         0.076       0.070
##   90 Percent confidence interval - upper         0.080       0.074
##   P-value H_0: RMSEA <= 0.050                    0.000       0.000
##   P-value H_0: RMSEA >= 0.080                    0.037       0.000
##                                                                   
##   Robust RMSEA                                               0.077
##   90 Percent confidence interval - lower                     0.075
##   90 Percent confidence interval - upper                     0.080
##   P-value H_0: Robust RMSEA <= 0.050                         0.000
##   P-value H_0: Robust RMSEA >= 0.080                         0.025
## 
## Standardized Root Mean Square Residual:
## 
##   SRMR                                           0.075       0.075
## 
## Parameter Estimates:
## 
##   Standard errors                             Sandwich
##   Information bread                           Observed
##   Observed information based on                Hessian
## 
## Latent Variables:
##                        Estimate  Std.Err  z-value  P(>|z|) ci.lower ci.upper
##   Agreeableness =~                                                          
##     A1                    0.484    0.035   13.947    0.000    0.416    0.552
##     A2                   -0.764    0.028  -27.062    0.000   -0.820   -0.709
##     A3                   -0.983    0.029  -34.326    0.000   -1.039   -0.926
##     A4                   -0.757    0.033  -22.614    0.000   -0.823   -0.692
##     A5                   -0.873    0.028  -31.503    0.000   -0.928   -0.819
##   Conscientiousness =~                                                      
##     C1                    0.680    0.034   20.030    0.000    0.614    0.747
##     C2                    0.781    0.035   22.326    0.000    0.712    0.849
##     C3                    0.705    0.030   23.655    0.000    0.646    0.763
##     C4                   -0.967    0.030  -32.184    0.000   -1.026   -0.908
##     C5                   -1.013    0.036  -28.006    0.000   -1.083   -0.942
##   Extraversion =~                                                           
##     E1                    0.920    0.035   26.422    0.000    0.852    0.988
##     E2                    1.128    0.031   36.139    0.000    1.066    1.189
##     E3                   -0.847    0.030  -27.791    0.000   -0.907   -0.788
##     E4                   -1.031    0.030  -33.850    0.000   -1.091   -0.972
##     E5                   -0.743    0.032  -23.492    0.000   -0.805   -0.681
##   Neuroticism =~                                                            
##     N1                    1.300    0.027   49.035    0.000    1.248    1.352
##     N2                    1.230    0.026   46.599    0.000    1.179    1.282
##     N3                    1.149    0.029   40.118    0.000    1.093    1.205
##     N4                    0.899    0.036   25.092    0.000    0.829    0.969
##     N5                    0.816    0.036   22.715    0.000    0.746    0.886
##   Openness =~                                                               
##     O1                    0.635    0.029   22.146    0.000    0.579    0.692
##     O2                   -0.648    0.046  -14.041    0.000   -0.739   -0.558
##     O3                    0.872    0.035   24.798    0.000    0.803    0.941
##     O4                    0.277    0.033    8.324    0.000    0.212    0.343
##     O5                   -0.610    0.039  -15.800    0.000   -0.685   -0.534
##    Std.lv  Std.all
##                   
##     0.484    0.344
##    -0.764   -0.648
##    -0.983   -0.749
##    -0.757   -0.510
##    -0.873   -0.687
##                   
##     0.680    0.551
##     0.781    0.592
##     0.705    0.546
##    -0.967   -0.702
##    -1.013   -0.620
##                   
##     0.920    0.564
##     1.128    0.699
##    -0.847   -0.627
##    -1.031   -0.703
##    -0.743   -0.553
##                   
##     1.300    0.825
##     1.230    0.803
##     1.149    0.721
##     0.899    0.573
##     0.816    0.503
##                   
##     0.635    0.564
##    -0.648   -0.418
##     0.872    0.724
##     0.277    0.233
##    -0.610   -0.461
## 
## Covariances:
##                        Estimate  Std.Err  z-value  P(>|z|) ci.lower ci.upper
##   Agreeableness ~~                                                          
##     Conscientisnss       -0.334    0.029  -11.446    0.000   -0.391   -0.277
##     Extraversion          0.683    0.022   31.135    0.000    0.640    0.726
##     Neuroticism           0.223    0.027    8.221    0.000    0.170    0.277
##     Openness             -0.303    0.033   -9.174    0.000   -0.368   -0.239
##   Conscientiousness ~~                                                      
##     Extraversion         -0.357    0.029  -12.341    0.000   -0.414   -0.301
##     Neuroticism          -0.283    0.030   -9.543    0.000   -0.341   -0.225
##     Openness              0.301    0.032    9.406    0.000    0.238    0.364
##   Extraversion ~~                                                           
##     Neuroticism           0.244    0.028    8.662    0.000    0.189    0.299
##     Openness             -0.453    0.035  -13.000    0.000   -0.521   -0.385
##   Neuroticism ~~                                                            
##     Openness             -0.112    0.030   -3.685    0.000   -0.172   -0.052
##    Std.lv  Std.all
##                   
##    -0.334   -0.334
##     0.683    0.683
##     0.223    0.223
##    -0.303   -0.303
##                   
##    -0.357   -0.357
##    -0.283   -0.283
##     0.301    0.301
##                   
##     0.244    0.244
##    -0.453   -0.453
##                   
##    -0.112   -0.112
## 
## Variances:
##                    Estimate  Std.Err  z-value  P(>|z|) ci.lower ci.upper
##    .A1                1.745    0.058   29.920    0.000    1.631    1.859
##    .A2                0.807    0.042   19.134    0.000    0.724    0.889
##    .A3                0.754    0.046   16.464    0.000    0.664    0.843
##    .A4                1.632    0.060   27.148    0.000    1.514    1.749
##    .A5                0.852    0.046   18.657    0.000    0.762    0.941
##    .C1                1.063    0.053   20.092    0.000    0.959    1.166
##    .C2                1.130    0.052   21.820    0.000    1.028    1.231
##    .C3                1.170    0.044   26.585    0.000    1.084    1.256
##    .C4                0.960    0.054   17.824    0.000    0.855    1.066
##    .C5                1.640    0.069   23.873    0.000    1.505    1.774
##    .E1                1.814    0.061   29.745    0.000    1.694    1.934
##    .E2                1.332    0.062   21.647    0.000    1.211    1.453
##    .E3                1.108    0.046   23.844    0.000    1.017    1.199
##    .E4                1.088    0.051   21.169    0.000    0.987    1.188
##    .E5                1.251    0.047   26.587    0.000    1.159    1.344
##    .N1                0.793    0.050   15.831    0.000    0.695    0.891
##    .N2                0.836    0.049   17.060    0.000    0.740    0.932
##    .N3                1.222    0.053   23.117    0.000    1.119    1.326
##    .N4                1.654    0.060   27.546    0.000    1.537    1.772
##    .N5                1.969    0.061   32.245    0.000    1.849    2.088
##    .O1                0.865    0.038   22.625    0.000    0.790    0.940
##    .O2                1.990    0.072   27.796    0.000    1.850    2.131
##    .O3                0.691    0.053   13.126    0.000    0.588    0.794
##    .O4                1.346    0.054   25.153    0.000    1.241    1.451
##    .O5                1.380    0.060   22.984    0.000    1.263    1.498
##     Agreeableness     1.000                               1.000    1.000
##     Conscientisnss    1.000                               1.000    1.000
##     Extraversion      1.000                               1.000    1.000
##     Neuroticism       1.000                               1.000    1.000
##     Openness          1.000                               1.000    1.000
##    Std.lv  Std.all
##     1.745    0.882
##     0.807    0.580
##     0.754    0.438
##     1.632    0.740
##     0.852    0.528
##     1.063    0.697
##     1.130    0.650
##     1.170    0.702
##     0.960    0.507
##     1.640    0.615
##     1.814    0.682
##     1.332    0.512
##     1.108    0.607
##     1.088    0.506
##     1.251    0.694
##     0.793    0.320
##     0.836    0.356
##     1.222    0.481
##     1.654    0.672
##     1.969    0.747
##     0.865    0.682
##     1.990    0.826
##     0.691    0.476
##     1.346    0.946
##     1.380    0.788
##     1.000    1.000
##     1.000    1.000
##     1.000    1.000
##     1.000    1.000
##     1.000    1.000
\end{verbatim}

\subsection{Reading the lavaan Output --- Complete Section-by-Section
Guide}\label{reading-the-lavaan-output-complete-section-by-section-guide}

The \texttt{summary()} output is long and contains many sections. Here
is what each one means and how to read it:

\begin{table}
\centering
\caption{\label{tab:lavaan-output-guide}Complete Section-by-Section Guide to lavaan CFA Output}
\centering
\begin{tabu} to \linewidth {>{\raggedright}X>{\raggedright}X}
\hline
Section & What It Means\\
\hline
\ttfamily{\textbf{lavaan version + Number of observations}} & Confirms the sample size. Always verify this matches your expected N after listwise deletion.\\
\hline
\ttfamily{\textbf{Estimator}} & MLR = Maximum Likelihood with robust SEs. Use MLR for non-normal data (ordinal items). Use ML for multivariate normal data.\\
\hline
\ttfamily{\textbf{Optimization method}} & NLMINB is the default optimizer. You do not need to change this.\\
\hline
\ttfamily{\textbf{Number of model parameters}} & The number of free parameters estimated. More parameters = more complex model.\\
\hline
\ttfamily{\textbf{Model Test User Model}} & The overall model chi-squared test\\
\hline
\ttfamily{\textbf{Test statistic (χ²)}} & How different the model-implied covariance matrix is from the observed one. Scaled by N. Almost always significant with large samples.\\
\hline
\ttfamily{\textbf{Degrees of freedom}} & Degrees of freedom = (p*(p+1)/2) − free parameters, where p = number of items. More df = more constrained model.\\
\hline
\ttfamily{\textbf{P-value}} & With N > 500, this is almost always < .05 even for good models. Focus on fit indices instead.\\
\hline
\ttfamily{\textbf{Model Test Baseline Model}} & The null model assumes all items are uncorrelated. CFI and TLI are computed relative to this null.\\
\hline
\ttfamily{\textbf{User Model vs. Baseline Model (CFI, TLI)}} & CFI: proportion of covariance explained vs. null. TLI: similar but penalises complexity. Both should be ≥ .95 for good fit.\\
\hline
\ttfamily{\textbf{Root Mean Square Error (RMSEA + CI)}} & RMSEA: average discrepancy per df. ≤ .06 = good, ≤ .08 = acceptable. The 90\% CI should not include values > .10.\\
\hline
\ttfamily{\textbf{Standardized Root Mean Square Residual (SRMR)}} & Average absolute difference between model-implied and observed correlations. ≤ .08 = good fit.\\
\hline
\ttfamily{\textbf{Parameter Estimates: Latent Variables (=\textasciitilde{})}} & This is the core output: estimated factor loadings\\
\hline
\ttfamily{\textbf{Estimate column}} & The unstandardised loading in the metric of the item (raw score units). Interpretation depends on scale.\\
\hline
\ttfamily{\textbf{Std.Err column}} & Standard error of the estimate. Smaller = more precise.\\
\hline
\ttfamily{\textbf{z-value column}} & Estimate / Std.Err. Used to compute p-value. |z| > 1.96 = p < .05.\\
\hline
\ttfamily{\textbf{P(>|z|) column}} & Two-tailed p-value testing H₀: this parameter = 0. Almost all loadings will be significant in a correctly specified model.\\
\hline
\ttfamily{\textbf{ci.lower / ci.upper columns}} & 95\% confidence interval. Does it include zero? If yes, the loading is not significant.\\
\hline
\ttfamily{\textbf{Std.lv column}} & Loading standardised with respect to latent variable variance only. Useful if std.lv=FALSE.\\
\hline
\ttfamily{\textbf{Std.all column}} & FULLY STANDARDISED LOADING: this is what you report in papers. Correlation between item and factor. Range: −1 to +1. > |.40| = acceptable; > |.60| = good; > |.80| = excellent.\\
\hline
\ttfamily{\textbf{Parameter Estimates: Covariances (\textasciitilde{}\textasciitilde{})}} & Factor covariances (or correlations if std.lv=TRUE). A covariance of .35 between Agreeableness and Conscientiousness means people high in one tend to be high in the other.\\
\hline
\ttfamily{\textbf{Parameter Estimates: Variances}} & Residual variance estimates for each item and each latent variable\\
\hline
\ttfamily{\textbf{Residual variances (for observed items)}} & Residual (unique) variance for each item — the portion NOT explained by its factor. 1 − Std.all² = residual proportion.\\
\hline
\ttfamily{\textbf{Variance of latent variables}} & If std.lv=TRUE, latent variances are fixed to 1.0 and will not appear as free parameters.\\
\hline
\end{tabu}
\end{table}

\subsection{Extracting and Displaying the Key
Results}\label{extracting-and-displaying-the-key-results}

\begin{Shaded}
\begin{Highlighting}[]
\CommentTok{\# Extract standardised loadings for clean reporting}
\NormalTok{loadings\_tbl }\OtherTok{\textless{}{-}} \FunctionTok{parameterEstimates}\NormalTok{(cfa\_fit, }\AttributeTok{standardized=}\ConstantTok{TRUE}\NormalTok{) }\SpecialCharTok{|\textgreater{}}
  \FunctionTok{filter}\NormalTok{(op }\SpecialCharTok{==} \StringTok{"=\textasciitilde{}"}\NormalTok{) }\SpecialCharTok{|\textgreater{}}      \CommentTok{\# keep only factor loading parameters}
\NormalTok{  dplyr}\SpecialCharTok{::}\FunctionTok{select}\NormalTok{(lhs, rhs, std.all, se, pvalue, ci.lower, ci.upper) }\SpecialCharTok{|\textgreater{}}
  \FunctionTok{mutate}\NormalTok{(}
    \AttributeTok{std.all  =} \FunctionTok{round}\NormalTok{(std.all, }\DecValTok{3}\NormalTok{),}
    \AttributeTok{se       =} \FunctionTok{round}\NormalTok{(se, }\DecValTok{3}\NormalTok{),}
    \AttributeTok{pvalue   =} \FunctionTok{ifelse}\NormalTok{(pvalue }\SpecialCharTok{\textless{}}\NormalTok{ .}\DecValTok{001}\NormalTok{, }\StringTok{"\textless{} .001"}\NormalTok{, }\FunctionTok{round}\NormalTok{(pvalue, }\DecValTok{3}\NormalTok{)),}
    \AttributeTok{ci.lower =} \FunctionTok{round}\NormalTok{(ci.lower, }\DecValTok{3}\NormalTok{),}
    \AttributeTok{ci.upper =} \FunctionTok{round}\NormalTok{(ci.upper, }\DecValTok{3}\NormalTok{),}
    \CommentTok{\# Communality from standardised loading: h2 = lambda\^{}2}
    \AttributeTok{h2       =} \FunctionTok{round}\NormalTok{(std.all}\SpecialCharTok{\^{}}\DecValTok{2}\NormalTok{, }\DecValTok{3}\NormalTok{),}
    \AttributeTok{Quality  =} \FunctionTok{case\_when}\NormalTok{(}
      \FunctionTok{abs}\NormalTok{(std.all) }\SpecialCharTok{\textgreater{}=} \FloatTok{0.70} \SpecialCharTok{\textasciitilde{}} \StringTok{"Strong"}\NormalTok{,}
      \FunctionTok{abs}\NormalTok{(std.all) }\SpecialCharTok{\textgreater{}=} \FloatTok{0.50} \SpecialCharTok{\textasciitilde{}} \StringTok{"Acceptable"}\NormalTok{,}
      \FunctionTok{abs}\NormalTok{(std.all) }\SpecialCharTok{\textgreater{}=} \FloatTok{0.30} \SpecialCharTok{\textasciitilde{}} \StringTok{"Marginal"}\NormalTok{,}
      \ConstantTok{TRUE}                 \SpecialCharTok{\textasciitilde{}} \StringTok{"Weak — review"}
\NormalTok{    )}
\NormalTok{  ) }\SpecialCharTok{|\textgreater{}}
  \FunctionTok{rename}\NormalTok{(}\AttributeTok{Factor=}\NormalTok{lhs, }\AttributeTok{Item=}\NormalTok{rhs, }\StringTok{\textasciigrave{}}\AttributeTok{λ (std.)}\StringTok{\textasciigrave{}}\OtherTok{=}\NormalTok{std.all, }\AttributeTok{SE=}\NormalTok{se, }\AttributeTok{p=}\NormalTok{pvalue,}
         \StringTok{\textasciigrave{}}\AttributeTok{CI low}\StringTok{\textasciigrave{}}\OtherTok{=}\NormalTok{ci.lower, }\StringTok{\textasciigrave{}}\AttributeTok{CI high}\StringTok{\textasciigrave{}}\OtherTok{=}\NormalTok{ci.upper, }\StringTok{\textasciigrave{}}\AttributeTok{h²}\StringTok{\textasciigrave{}}\OtherTok{=}\NormalTok{h2)}

\FunctionTok{kbl}\NormalTok{(loadings\_tbl,}
    \AttributeTok{caption=}\StringTok{"CFA Standardised Factor Loadings — Big Five Model"}\NormalTok{) }\SpecialCharTok{|\textgreater{}}
  \FunctionTok{kable\_styling}\NormalTok{(}\AttributeTok{bootstrap\_options=}\FunctionTok{c}\NormalTok{(}\StringTok{"striped"}\NormalTok{,}\StringTok{"hover"}\NormalTok{,}\StringTok{"condensed"}\NormalTok{)) }\SpecialCharTok{|\textgreater{}}
  \FunctionTok{group\_rows}\NormalTok{(}\StringTok{"Agreeableness"}\NormalTok{, }\DecValTok{1}\NormalTok{, }\DecValTok{5}\NormalTok{) }\SpecialCharTok{|\textgreater{}}
  \FunctionTok{group\_rows}\NormalTok{(}\StringTok{"Conscientiousness"}\NormalTok{, }\DecValTok{6}\NormalTok{, }\DecValTok{10}\NormalTok{) }\SpecialCharTok{|\textgreater{}}
  \FunctionTok{group\_rows}\NormalTok{(}\StringTok{"Extraversion"}\NormalTok{, }\DecValTok{11}\NormalTok{, }\DecValTok{15}\NormalTok{) }\SpecialCharTok{|\textgreater{}}
  \FunctionTok{group\_rows}\NormalTok{(}\StringTok{"Neuroticism"}\NormalTok{, }\DecValTok{16}\NormalTok{, }\DecValTok{20}\NormalTok{) }\SpecialCharTok{|\textgreater{}}
  \FunctionTok{group\_rows}\NormalTok{(}\StringTok{"Openness"}\NormalTok{, }\DecValTok{21}\NormalTok{, }\DecValTok{25}\NormalTok{) }\SpecialCharTok{|\textgreater{}}
  \FunctionTok{column\_spec}\NormalTok{(}\DecValTok{8}\NormalTok{, }\AttributeTok{color=}\FunctionTok{ifelse}\NormalTok{(loadings\_tbl}\SpecialCharTok{$}\NormalTok{Quality }\SpecialCharTok{\%in\%} \FunctionTok{c}\NormalTok{(}\StringTok{"Weak — review"}\NormalTok{,}\StringTok{"Marginal"}\NormalTok{),}
                               \StringTok{"\#C0392B"}\NormalTok{,}\StringTok{"\#27AE60"}\NormalTok{))}
\end{Highlighting}
\end{Shaded}

\begin{table}
\centering
\caption{\label{tab:cfa-loadings-clean}CFA Standardised Factor Loadings — Big Five Model}
\centering
\begin{tabular}[t]{l|l|r|r|l|r|r|>{}r|l}
\hline
Factor & Item & λ (std.) & SE & p & CI low & CI high & h² & Quality\\
\hline
\multicolumn{9}{l}{\textbf{Agreeableness}}\\
\hline
\hspace{1em}Agreeableness & A1 & 0.344 & 0.035 & < .001 & 0.416 & 0.552 & \textcolor[HTML]{C0392B}{0.118} & Marginal\\
\hline
\hspace{1em}Agreeableness & A2 & -0.648 & 0.028 & < .001 & -0.820 & -0.709 & \textcolor[HTML]{27AE60}{0.420} & Acceptable\\
\hline
\hspace{1em}Agreeableness & A3 & -0.749 & 0.029 & < .001 & -1.039 & -0.926 & \textcolor[HTML]{27AE60}{0.561} & Strong\\
\hline
\hspace{1em}Agreeableness & A4 & -0.510 & 0.033 & < .001 & -0.823 & -0.692 & \textcolor[HTML]{27AE60}{0.260} & Acceptable\\
\hline
\hspace{1em}Agreeableness & A5 & -0.687 & 0.028 & < .001 & -0.928 & -0.819 & \textcolor[HTML]{27AE60}{0.472} & Acceptable\\
\hline
\multicolumn{9}{l}{\textbf{Conscientiousness}}\\
\hline
\hspace{1em}Conscientiousness & C1 & 0.551 & 0.034 & < .001 & 0.614 & 0.747 & \textcolor[HTML]{27AE60}{0.304} & Acceptable\\
\hline
\hspace{1em}Conscientiousness & C2 & 0.592 & 0.035 & < .001 & 0.712 & 0.849 & \textcolor[HTML]{27AE60}{0.350} & Acceptable\\
\hline
\hspace{1em}Conscientiousness & C3 & 0.546 & 0.030 & < .001 & 0.646 & 0.763 & \textcolor[HTML]{27AE60}{0.298} & Acceptable\\
\hline
\hspace{1em}Conscientiousness & C4 & -0.702 & 0.030 & < .001 & -1.026 & -0.908 & \textcolor[HTML]{27AE60}{0.493} & Strong\\
\hline
\hspace{1em}Conscientiousness & C5 & -0.620 & 0.036 & < .001 & -1.083 & -0.942 & \textcolor[HTML]{27AE60}{0.384} & Acceptable\\
\hline
\multicolumn{9}{l}{\textbf{Extraversion}}\\
\hline
\hspace{1em}Extraversion & E1 & 0.564 & 0.035 & < .001 & 0.852 & 0.988 & \textcolor[HTML]{27AE60}{0.318} & Acceptable\\
\hline
\hspace{1em}Extraversion & E2 & 0.699 & 0.031 & < .001 & 1.066 & 1.189 & \textcolor[HTML]{27AE60}{0.489} & Acceptable\\
\hline
\hspace{1em}Extraversion & E3 & -0.627 & 0.030 & < .001 & -0.907 & -0.788 & \textcolor[HTML]{27AE60}{0.393} & Acceptable\\
\hline
\hspace{1em}Extraversion & E4 & -0.703 & 0.030 & < .001 & -1.091 & -0.972 & \textcolor[HTML]{27AE60}{0.494} & Strong\\
\hline
\hspace{1em}Extraversion & E5 & -0.553 & 0.032 & < .001 & -0.805 & -0.681 & \textcolor[HTML]{27AE60}{0.306} & Acceptable\\
\hline
\multicolumn{9}{l}{\textbf{Neuroticism}}\\
\hline
\hspace{1em}Neuroticism & N1 & 0.825 & 0.027 & < .001 & 1.248 & 1.352 & \textcolor[HTML]{27AE60}{0.681} & Strong\\
\hline
\hspace{1em}Neuroticism & N2 & 0.803 & 0.026 & < .001 & 1.179 & 1.282 & \textcolor[HTML]{27AE60}{0.645} & Strong\\
\hline
\hspace{1em}Neuroticism & N3 & 0.721 & 0.029 & < .001 & 1.093 & 1.205 & \textcolor[HTML]{27AE60}{0.520} & Strong\\
\hline
\hspace{1em}Neuroticism & N4 & 0.573 & 0.036 & < .001 & 0.829 & 0.969 & \textcolor[HTML]{27AE60}{0.328} & Acceptable\\
\hline
\hspace{1em}Neuroticism & N5 & 0.503 & 0.036 & < .001 & 0.746 & 0.886 & \textcolor[HTML]{27AE60}{0.253} & Acceptable\\
\hline
\multicolumn{9}{l}{\textbf{Openness}}\\
\hline
\hspace{1em}Openness & O1 & 0.564 & 0.029 & < .001 & 0.579 & 0.692 & \textcolor[HTML]{27AE60}{0.318} & Acceptable\\
\hline
\hspace{1em}Openness & O2 & -0.418 & 0.046 & < .001 & -0.739 & -0.558 & \textcolor[HTML]{C0392B}{0.175} & Marginal\\
\hline
\hspace{1em}Openness & O3 & 0.724 & 0.035 & < .001 & 0.803 & 0.941 & \textcolor[HTML]{27AE60}{0.524} & Strong\\
\hline
\hspace{1em}Openness & O4 & 0.233 & 0.033 & < .001 & 0.212 & 0.343 & \textcolor[HTML]{C0392B}{0.054} & Weak — review\\
\hline
\hspace{1em}Openness & O5 & -0.461 & 0.039 & < .001 & -0.685 & -0.534 & \textcolor[HTML]{C0392B}{0.213} & Marginal\\
\hline
\end{tabular}
\end{table}

\textbf{Reading this table:}

\begin{itemize}
\tightlist
\item
  \textbf{λ (std.):} The standardised loading. This is the number you
  report in your methods section. A loading of .65 means the item
  correlates .65 with its factor.
\item
  \textbf{h²:} Communality = λ². An h² of .42 means 42\% of this item's
  variance is shared with its factor; 58\% is unique. Items with h²
  \textless{} .20 are barely measuring the factor.
\item
  \textbf{Quality column:} Benchmark guidance. A scale with mostly
  ``Acceptable'' and ``Strong'' items is well-specified. Several
  ``Weak'' items suggest the items need revision or the model needs
  rethinking.
\item
  \textbf{p column:} Almost every loading will be significant --- this
  is just telling you the loading is not zero. The magnitude (λ) and
  communality (h²) are more informative than the p-value here.
\end{itemize}

\subsection{Evaluating Model Fit}\label{evaluating-model-fit}

\begin{Shaded}
\begin{Highlighting}[]
\NormalTok{fit\_indices }\OtherTok{\textless{}{-}} \FunctionTok{fitMeasures}\NormalTok{(cfa\_fit,}
  \FunctionTok{c}\NormalTok{(}\StringTok{"chisq"}\NormalTok{,}\StringTok{"df"}\NormalTok{,}\StringTok{"pvalue"}\NormalTok{,}\StringTok{"cfi"}\NormalTok{,}\StringTok{"tli"}\NormalTok{,}\StringTok{"rmsea"}\NormalTok{,}\StringTok{"rmsea.ci.lower"}\NormalTok{,}
    \StringTok{"rmsea.ci.upper"}\NormalTok{,}\StringTok{"rmsea.pvalue"}\NormalTok{,}\StringTok{"srmr"}\NormalTok{,}\StringTok{"aic"}\NormalTok{,}\StringTok{"bic"}\NormalTok{)) }\SpecialCharTok{|\textgreater{}}
  \FunctionTok{round}\NormalTok{(}\DecValTok{3}\NormalTok{)}

\FunctionTok{tibble}\NormalTok{(}
  \AttributeTok{Index   =} \FunctionTok{c}\NormalTok{(}\StringTok{"χ² (chi{-}squared)"}\NormalTok{,}\StringTok{"df"}\NormalTok{,}\StringTok{"p(χ²)"}\NormalTok{,}\StringTok{"CFI"}\NormalTok{,}\StringTok{"TLI"}\NormalTok{,}
               \StringTok{"RMSEA"}\NormalTok{,}\StringTok{"RMSEA 95\% CI lower"}\NormalTok{,}\StringTok{"RMSEA 95\% CI upper"}\NormalTok{,}
               \StringTok{"p(RMSEA ≤ .05)"}\NormalTok{,}\StringTok{"SRMR"}\NormalTok{,}\StringTok{"AIC"}\NormalTok{,}\StringTok{"BIC"}\NormalTok{),}
  \AttributeTok{Value   =}\NormalTok{ fit\_indices,}
  \AttributeTok{Benchmark =} \FunctionTok{c}\NormalTok{(}\StringTok{"—"}\NormalTok{,}\StringTok{"—"}\NormalTok{,}\StringTok{"\textgreater{} .05 for exact fit"}\NormalTok{,}
                \StringTok{"≥ .95 good, ≥ .90 acceptable"}\NormalTok{,}
                \StringTok{"≥ .95 good, ≥ .90 acceptable"}\NormalTok{,}
                \StringTok{"≤ .06 good, ≤ .08 acceptable"}\NormalTok{,}
                \StringTok{"—"}\NormalTok{,}\StringTok{"—"}\NormalTok{,}\StringTok{"\textgreater{}.05 = close fit"}\NormalTok{,}
                \StringTok{"≤ .08 good, ≤ .10 acceptable"}\NormalTok{,}
                \StringTok{"Lower = better"}\NormalTok{,}\StringTok{"Lower = better"}\NormalTok{),}
  \AttributeTok{Verdict =} \FunctionTok{c}\NormalTok{(}
    \StringTok{"Ignore significance with N \textgreater{} 300"}\NormalTok{,}
    \StringTok{"More df = more constrained (parsimonious) model"}\NormalTok{,}
    \ControlFlowTok{if}\NormalTok{(fit\_indices[}\StringTok{"pvalue"}\NormalTok{] }\SpecialCharTok{\textgreater{}}\NormalTok{ .}\DecValTok{05}\NormalTok{) }\StringTok{"✓ Not significant — exact fit"} \ControlFlowTok{else} \StringTok{"✗ Significant (expected at large N)"}\NormalTok{,}
    \ControlFlowTok{if}\NormalTok{(fit\_indices[}\StringTok{"cfi"}\NormalTok{] }\SpecialCharTok{\textgreater{}=}\NormalTok{ .}\DecValTok{95}\NormalTok{) }\StringTok{"✓ Good"} \ControlFlowTok{else} \ControlFlowTok{if}\NormalTok{(fit\_indices[}\StringTok{"cfi"}\NormalTok{] }\SpecialCharTok{\textgreater{}=}\NormalTok{ .}\DecValTok{90}\NormalTok{) }\StringTok{"\textasciitilde{} Acceptable"} \ControlFlowTok{else} \StringTok{"✗ Poor"}\NormalTok{,}
    \ControlFlowTok{if}\NormalTok{(fit\_indices[}\StringTok{"tli"}\NormalTok{] }\SpecialCharTok{\textgreater{}=}\NormalTok{ .}\DecValTok{95}\NormalTok{) }\StringTok{"✓ Good"} \ControlFlowTok{else} \ControlFlowTok{if}\NormalTok{(fit\_indices[}\StringTok{"tli"}\NormalTok{] }\SpecialCharTok{\textgreater{}=}\NormalTok{ .}\DecValTok{90}\NormalTok{) }\StringTok{"\textasciitilde{} Acceptable"} \ControlFlowTok{else} \StringTok{"✗ Poor"}\NormalTok{,}
    \ControlFlowTok{if}\NormalTok{(fit\_indices[}\StringTok{"rmsea"}\NormalTok{] }\SpecialCharTok{\textless{}=}\NormalTok{ .}\DecValTok{06}\NormalTok{) }\StringTok{"✓ Good"} \ControlFlowTok{else} \ControlFlowTok{if}\NormalTok{(fit\_indices[}\StringTok{"rmsea"}\NormalTok{] }\SpecialCharTok{\textless{}=}\NormalTok{ .}\DecValTok{08}\NormalTok{) }\StringTok{"\textasciitilde{} Acceptable"} \ControlFlowTok{else} \StringTok{"✗ Poor"}\NormalTok{,}
    \StringTok{"—"}\NormalTok{,}\StringTok{"—"}\NormalTok{,}
    \ControlFlowTok{if}\NormalTok{(fit\_indices[}\StringTok{"rmsea.pvalue"}\NormalTok{] }\SpecialCharTok{\textgreater{}}\NormalTok{ .}\DecValTok{05}\NormalTok{) }\StringTok{"✓ Close fit confirmed"} \ControlFlowTok{else} \StringTok{"\textasciitilde{} Not close fit"}\NormalTok{,}
    \ControlFlowTok{if}\NormalTok{(fit\_indices[}\StringTok{"srmr"}\NormalTok{] }\SpecialCharTok{\textless{}=}\NormalTok{ .}\DecValTok{08}\NormalTok{) }\StringTok{"✓ Good"} \ControlFlowTok{else} \ControlFlowTok{if}\NormalTok{(fit\_indices[}\StringTok{"srmr"}\NormalTok{] }\SpecialCharTok{\textless{}=}\NormalTok{ .}\DecValTok{10}\NormalTok{) }\StringTok{"\textasciitilde{} Acceptable"} \ControlFlowTok{else} \StringTok{"✗ Poor"}\NormalTok{,}
    \StringTok{"Use for model comparison only"}\NormalTok{,}\StringTok{"Use for model comparison only"}
\NormalTok{  )}
\NormalTok{) }\SpecialCharTok{|\textgreater{}}
  \FunctionTok{kbl}\NormalTok{(}\AttributeTok{caption=}\StringTok{"CFA Fit Index Summary — Big Five Model"}\NormalTok{,}
      \AttributeTok{col.names=}\FunctionTok{c}\NormalTok{(}\StringTok{"Index"}\NormalTok{,}\StringTok{"Value"}\NormalTok{,}\StringTok{"Benchmark"}\NormalTok{,}\StringTok{"Verdict"}\NormalTok{)) }\SpecialCharTok{|\textgreater{}}
  \FunctionTok{kable\_styling}\NormalTok{(}\AttributeTok{bootstrap\_options=}\FunctionTok{c}\NormalTok{(}\StringTok{"striped"}\NormalTok{,}\StringTok{"hover"}\NormalTok{)) }\SpecialCharTok{|\textgreater{}}
  \FunctionTok{row\_spec}\NormalTok{(}\FunctionTok{c}\NormalTok{(}\DecValTok{4}\NormalTok{,}\DecValTok{5}\NormalTok{), }\AttributeTok{background=}\StringTok{"\#EBF5FB"}\NormalTok{) }\SpecialCharTok{|\textgreater{}}
  \FunctionTok{row\_spec}\NormalTok{(}\DecValTok{6}\NormalTok{, }\AttributeTok{background=}\StringTok{"\#EAFAF1"}\NormalTok{) }\SpecialCharTok{|\textgreater{}}
  \FunctionTok{row\_spec}\NormalTok{(}\DecValTok{10}\NormalTok{, }\AttributeTok{background=}\StringTok{"\#FEF9E7"}\NormalTok{) }\SpecialCharTok{|\textgreater{}}
  \FunctionTok{footnote}\NormalTok{(}\AttributeTok{general=}\FunctionTok{paste}\NormalTok{(}
    \StringTok{"CFI and TLI compare your model to a null model where all items are uncorrelated."}\NormalTok{,}
    \StringTok{"}\SpecialCharTok{\textbackslash{}n}\StringTok{RMSEA measures the average discrepancy per degree of freedom."}\NormalTok{,}
    \StringTok{"}\SpecialCharTok{\textbackslash{}n}\StringTok{SRMR is the average absolute difference between model{-}implied and observed correlations."}
\NormalTok{  ))}
\end{Highlighting}
\end{Shaded}

\begin{table}
\centering
\caption{\label{tab:fit-indices-table}CFA Fit Index Summary — Big Five Model}
\centering
\begin{tabular}[t]{l|r|l|l}
\hline
Index & Value & Benchmark & Verdict\\
\hline
χ² (chi-squared) & 4165.467 & — & Ignore significance with N > 300\\
\hline
df & 265.000 & — & More df = more constrained (parsimonious) model\\
\hline
p(χ²) & 0.000 & > .05 for exact fit & ✗ Significant (expected at large N)\\
\hline
\cellcolor[HTML]{EBF5FB}{CFI} & \cellcolor[HTML]{EBF5FB}{0.782} & \cellcolor[HTML]{EBF5FB}{≥ .95 good, ≥ .90 acceptable} & \cellcolor[HTML]{EBF5FB}{✗ Poor}\\
\hline
\cellcolor[HTML]{EBF5FB}{TLI} & \cellcolor[HTML]{EBF5FB}{0.754} & \cellcolor[HTML]{EBF5FB}{≥ .95 good, ≥ .90 acceptable} & \cellcolor[HTML]{EBF5FB}{✗ Poor}\\
\hline
\cellcolor[HTML]{EAFAF1}{RMSEA} & \cellcolor[HTML]{EAFAF1}{0.078} & \cellcolor[HTML]{EAFAF1}{≤ .06 good, ≤ .08 acceptable} & \cellcolor[HTML]{EAFAF1}{\textasciitilde{} Acceptable}\\
\hline
RMSEA 95\% CI lower & 0.076 & — & —\\
\hline
RMSEA 95\% CI upper & 0.080 & — & —\\
\hline
p(RMSEA ≤ .05) & 0.000 & >.05 = close fit & \textasciitilde{} Not close fit\\
\hline
\cellcolor[HTML]{FEF9E7}{SRMR} & \cellcolor[HTML]{FEF9E7}{0.075} & \cellcolor[HTML]{FEF9E7}{≤ .08 good, ≤ .10 acceptable} & \cellcolor[HTML]{FEF9E7}{✓ Good}\\
\hline
AIC & 199800.476 & Lower = better & Use for model comparison only\\
\hline
BIC & 200148.363 & Lower = better & Use for model comparison only\\
\hline
\multicolumn{4}{l}{\rule{0pt}{1em}\textit{Note: }}\\
\multicolumn{4}{l}{\rule{0pt}{1em}makecell[l]{CFI and TLI compare your model to a null model where all items are uncorrelated. \RMSEA measures the average discrepancy per degree of freedom. \SRMR is the average absolute difference between model-implied and observed correlations.}}\\
\end{tabular}
\end{table}

\subsection{Factor Correlations}\label{factor-correlations}

\begin{Shaded}
\begin{Highlighting}[]
\CommentTok{\# Extract factor covariances/correlations from the model}
\CommentTok{\# With std.lv=TRUE, these ARE correlations (variances fixed to 1)}
\NormalTok{factor\_cors }\OtherTok{\textless{}{-}} \FunctionTok{parameterEstimates}\NormalTok{(cfa\_fit, }\AttributeTok{standardized=}\ConstantTok{TRUE}\NormalTok{) }\SpecialCharTok{|\textgreater{}}
  \FunctionTok{filter}\NormalTok{(op }\SpecialCharTok{==} \StringTok{"\textasciitilde{}\textasciitilde{}"}\NormalTok{,}
\NormalTok{         lhs }\SpecialCharTok{!=}\NormalTok{ rhs,}
\NormalTok{         lhs }\SpecialCharTok{\%in\%} \FunctionTok{c}\NormalTok{(}\StringTok{"Agreeableness"}\NormalTok{,}\StringTok{"Conscientiousness"}\NormalTok{,}\StringTok{"Extraversion"}\NormalTok{,}
                     \StringTok{"Neuroticism"}\NormalTok{,}\StringTok{"Openness"}\NormalTok{)) }\SpecialCharTok{|\textgreater{}}
\NormalTok{  dplyr}\SpecialCharTok{::}\FunctionTok{select}\NormalTok{(lhs, rhs, std.all, pvalue) }\SpecialCharTok{|\textgreater{}}
  \FunctionTok{mutate}\NormalTok{(}\AttributeTok{std.all =} \FunctionTok{round}\NormalTok{(std.all, }\DecValTok{3}\NormalTok{),}
         \AttributeTok{pvalue  =} \FunctionTok{round}\NormalTok{(pvalue, }\DecValTok{4}\NormalTok{),}
         \AttributeTok{sig     =} \FunctionTok{ifelse}\NormalTok{(pvalue }\SpecialCharTok{\textless{}}\NormalTok{ .}\DecValTok{05}\NormalTok{, }\StringTok{"*"}\NormalTok{, }\StringTok{""}\NormalTok{))}

\FunctionTok{cat}\NormalTok{(}\StringTok{"=== Factor Correlations (from CFA) ===}\SpecialCharTok{\textbackslash{}n}\StringTok{"}\NormalTok{)}
\end{Highlighting}
\end{Shaded}

\begin{verbatim}
## === Factor Correlations (from CFA) ===
\end{verbatim}

\begin{Shaded}
\begin{Highlighting}[]
\FunctionTok{cat}\NormalTok{(}\StringTok{"(These are the correlations between latent factors,}\SpecialCharTok{\textbackslash{}n}\StringTok{"}\NormalTok{)}
\end{Highlighting}
\end{Shaded}

\begin{verbatim}
## (These are the correlations between latent factors,
\end{verbatim}

\begin{Shaded}
\begin{Highlighting}[]
\FunctionTok{cat}\NormalTok{(}\StringTok{"corrected for measurement error in the items)}\SpecialCharTok{\textbackslash{}n\textbackslash{}n}\StringTok{"}\NormalTok{)}
\end{Highlighting}
\end{Shaded}

\begin{verbatim}
## corrected for measurement error in the items)
\end{verbatim}

\begin{Shaded}
\begin{Highlighting}[]
\NormalTok{factor\_cors }\SpecialCharTok{|\textgreater{}}
\NormalTok{  dplyr}\SpecialCharTok{::}\FunctionTok{select}\NormalTok{(}\AttributeTok{Factor1=}\NormalTok{lhs, }\AttributeTok{Factor2=}\NormalTok{rhs, }\StringTok{\textasciigrave{}}\AttributeTok{r (latent)}\StringTok{\textasciigrave{}}\OtherTok{=}\NormalTok{std.all, }\AttributeTok{p=}\NormalTok{pvalue) }\SpecialCharTok{|\textgreater{}}
  \FunctionTok{kbl}\NormalTok{() }\SpecialCharTok{|\textgreater{}}
  \FunctionTok{kable\_styling}\NormalTok{(}\AttributeTok{bootstrap\_options=}\FunctionTok{c}\NormalTok{(}\StringTok{"striped"}\NormalTok{,}\StringTok{"hover"}\NormalTok{,}\StringTok{"condensed"}\NormalTok{), }\AttributeTok{full\_width=}\ConstantTok{FALSE}\NormalTok{) }\SpecialCharTok{|\textgreater{}}
  \FunctionTok{footnote}\NormalTok{(}\AttributeTok{general=}\FunctionTok{paste}\NormalTok{(}
    \StringTok{"Note: latent factor correlations are disattenuated — they correct for measurement error."}\NormalTok{,}
    \StringTok{"}\SpecialCharTok{\textbackslash{}n}\StringTok{They will typically be larger (in absolute value) than observed scale{-}score correlations."}\NormalTok{,}
    \StringTok{"}\SpecialCharTok{\textbackslash{}n}\StringTok{Negative correlation with Neuroticism is expected (neurotic people tend to be less agreeable, etc.)"}
\NormalTok{  ))}
\end{Highlighting}
\end{Shaded}

\begin{table}
\centering
\begin{tabular}[t]{l|l|r|r}
\hline
Factor1 & Factor2 & r (latent) & p\\
\hline
Agreeableness & Conscientiousness & -0.334 & 0e+00\\
\hline
Agreeableness & Extraversion & 0.683 & 0e+00\\
\hline
Agreeableness & Neuroticism & 0.223 & 0e+00\\
\hline
Agreeableness & Openness & -0.303 & 0e+00\\
\hline
Conscientiousness & Extraversion & -0.357 & 0e+00\\
\hline
Conscientiousness & Neuroticism & -0.283 & 0e+00\\
\hline
Conscientiousness & Openness & 0.301 & 0e+00\\
\hline
Extraversion & Neuroticism & 0.244 & 0e+00\\
\hline
Extraversion & Openness & -0.453 & 0e+00\\
\hline
Neuroticism & Openness & -0.112 & 2e-04\\
\hline
\multicolumn{4}{l}{\rule{0pt}{1em}\textit{Note: }}\\
\multicolumn{4}{l}{\rule{0pt}{1em}makecell[l]{Note: latent factor correlations are disattenuated — they correct for measurement error. \They will typically be larger (in absolute value) than observed scale-score correlations. \Negative correlation with Neuroticism is expected (neurotic people tend to be less agreeable, etc.)}}\\
\end{tabular}
\end{table}

\subsection{Path Diagram}\label{path-diagram}

\begin{Shaded}
\begin{Highlighting}[]
\FunctionTok{semPaths}\NormalTok{(}
\NormalTok{  cfa\_fit,}
  \AttributeTok{whatLabels =} \StringTok{"std"}\NormalTok{,        }\CommentTok{\# show standardised loadings on arrows}
  \AttributeTok{layout     =} \StringTok{"tree2"}\NormalTok{,      }\CommentTok{\# factors at top, items below}
  \AttributeTok{style      =} \StringTok{"lisrel"}\NormalTok{,     }\CommentTok{\# traditional box{-}and{-}arrow diagram}
  \AttributeTok{edge.label.cex =} \FloatTok{0.60}\NormalTok{,}
  \AttributeTok{residuals  =} \ConstantTok{TRUE}\NormalTok{,         }\CommentTok{\# show residual variances as loops}
  \AttributeTok{intercepts =} \ConstantTok{FALSE}\NormalTok{,        }\CommentTok{\# omit intercepts for cleanliness}
  \AttributeTok{nCharNodes =} \DecValTok{0}\NormalTok{,            }\CommentTok{\# show full variable names (no truncation)}
  \AttributeTok{mar        =} \FunctionTok{c}\NormalTok{(}\DecValTok{1}\NormalTok{, }\DecValTok{1}\NormalTok{, }\DecValTok{2}\NormalTok{, }\DecValTok{1}\NormalTok{)}
\NormalTok{)}
\FunctionTok{title}\NormalTok{(}\StringTok{"CFA Path Diagram — Big Five Personality (Standardised loadings)"}\NormalTok{, }\AttributeTok{cex.main=}\FloatTok{0.9}\NormalTok{)}
\end{Highlighting}
\end{Shaded}

\begin{center}\includegraphics{module_03_efa_cfa_files/figure-latex/path-diagram-1} \end{center}

\textbf{Reading the path diagram:}

\begin{itemize}
\tightlist
\item
  \textbf{Ovals (or circles):} Latent variables (factors)
\item
  \textbf{Rectangles (or squares):} Observed items
\item
  \textbf{Single-headed arrows (→):} Factor loadings (from factor to
  item)
\item
  \textbf{Numbers on arrows:} Standardised factor loadings
\item
  \textbf{Small curved arrows (loops):} Residual variances --- the
  variance in each item \emph{not} explained by its factor
\item
  \textbf{Double-headed arrows (↔) between ovals:} Factor correlations
\end{itemize}

Strong loadings (\textgreater{} .60) produce thick arrows. Weak loadings
(\textless{} .30) produce thin, barely visible arrows. If you see many
thin arrows in a real analysis, those items should be reconsidered.

\begin{center}\rule{0.5\linewidth}{0.5pt}\end{center}

\section{Reliability Estimation}\label{reliability-estimation}

\subsection{Cronbach's Alpha --- The Standard (But Imperfect)
Coefficient}\label{cronbachs-alpha-the-standard-but-imperfect-coefficient}

Cronbach's alpha (α) estimates the proportion of scale variance
attributable to the underlying construct. Values \textgreater{} .70 are
conventionally ``acceptable'' for research; \textgreater{} .80 for
applied settings; \textgreater{} .90 for clinical use.

\textbf{Alpha's hidden assumption:} All items have \emph{equal} factor
loadings (tau-equivalence). When loadings are unequal --- as they almost
always are --- alpha underestimates true reliability.

\begin{Shaded}
\begin{Highlighting}[]
\CommentTok{\# Compute alpha for the Extraversion subscale (E1–E5)}
\NormalTok{extra\_items }\OtherTok{\textless{}{-}}\NormalTok{ bfi\_complete[, }\FunctionTok{paste0}\NormalTok{(}\StringTok{"E"}\NormalTok{, }\DecValTok{1}\SpecialCharTok{:}\DecValTok{5}\NormalTok{)]}

\NormalTok{alpha\_res }\OtherTok{\textless{}{-}}\NormalTok{ psych}\SpecialCharTok{::}\FunctionTok{alpha}\NormalTok{(extra\_items)}
\end{Highlighting}
\end{Shaded}

\begin{verbatim}
## Some items ( E1 E2 ) were negatively correlated with the first principal component and 
## probably should be reversed.  
## To do this, run the function again with the 'check.keys=TRUE' option
\end{verbatim}

\begin{Shaded}
\begin{Highlighting}[]
\FunctionTok{cat}\NormalTok{(}\StringTok{"=== Reliability Analysis: Extraversion (E1–E5) ===}\SpecialCharTok{\textbackslash{}n\textbackslash{}n}\StringTok{"}\NormalTok{)}
\end{Highlighting}
\end{Shaded}

\begin{verbatim}
## === Reliability Analysis: Extraversion (E1–E5) ===
\end{verbatim}

\begin{Shaded}
\begin{Highlighting}[]
\FunctionTok{cat}\NormalTok{(}\StringTok{"Cronbach\textquotesingle{}s alpha (raw):     "}\NormalTok{, }\FunctionTok{round}\NormalTok{(alpha\_res}\SpecialCharTok{$}\NormalTok{total}\SpecialCharTok{$}\NormalTok{raw\_alpha, }\DecValTok{3}\NormalTok{), }\StringTok{"}\SpecialCharTok{\textbackslash{}n}\StringTok{"}\NormalTok{)}
\end{Highlighting}
\end{Shaded}

\begin{verbatim}
## Cronbach's alpha (raw):      -0.655
\end{verbatim}

\begin{Shaded}
\begin{Highlighting}[]
\FunctionTok{cat}\NormalTok{(}\StringTok{"Cronbach\textquotesingle{}s alpha (std.):    "}\NormalTok{, }\FunctionTok{round}\NormalTok{(alpha\_res}\SpecialCharTok{$}\NormalTok{total}\SpecialCharTok{$}\NormalTok{std.alpha, }\DecValTok{3}\NormalTok{), }\StringTok{"}\SpecialCharTok{\textbackslash{}n}\StringTok{"}\NormalTok{)}
\end{Highlighting}
\end{Shaded}

\begin{verbatim}
## Cronbach's alpha (std.):     -0.531
\end{verbatim}

\begin{Shaded}
\begin{Highlighting}[]
\FunctionTok{cat}\NormalTok{(}\StringTok{"Average inter{-}item r:       "}\NormalTok{, }\FunctionTok{round}\NormalTok{(alpha\_res}\SpecialCharTok{$}\NormalTok{total}\SpecialCharTok{$}\NormalTok{average\_r, }\DecValTok{3}\NormalTok{), }\StringTok{"}\SpecialCharTok{\textbackslash{}n}\StringTok{"}\NormalTok{)}
\end{Highlighting}
\end{Shaded}

\begin{verbatim}
## Average inter-item r:        -0.075
\end{verbatim}

\begin{Shaded}
\begin{Highlighting}[]
\FunctionTok{cat}\NormalTok{(}\StringTok{"Guttman\textquotesingle{}s lambda 6 (G6):    "}\NormalTok{, }\FunctionTok{round}\NormalTok{(alpha\_res}\SpecialCharTok{$}\NormalTok{total[[}\StringTok{"G6(smc)"}\NormalTok{]], }\DecValTok{3}\NormalTok{), }\StringTok{"}\SpecialCharTok{\textbackslash{}n}\StringTok{"}\NormalTok{)}
\end{Highlighting}
\end{Shaded}

\begin{verbatim}
## Guttman's lambda 6 (G6):     0.011
\end{verbatim}

\begin{Shaded}
\begin{Highlighting}[]
\FunctionTok{cat}\NormalTok{(}\StringTok{"95\% CI:                      ["}\NormalTok{,}
    \FunctionTok{round}\NormalTok{(alpha\_res}\SpecialCharTok{$}\NormalTok{total}\SpecialCharTok{$}\NormalTok{raw\_alpha }\SpecialCharTok{{-}} \FloatTok{1.96}\SpecialCharTok{*}\NormalTok{alpha\_res}\SpecialCharTok{$}\NormalTok{total}\SpecialCharTok{$}\NormalTok{ase, }\DecValTok{3}\NormalTok{), }\StringTok{","}\NormalTok{,}
    \FunctionTok{round}\NormalTok{(alpha\_res}\SpecialCharTok{$}\NormalTok{total}\SpecialCharTok{$}\NormalTok{raw\_alpha }\SpecialCharTok{+} \FloatTok{1.96}\SpecialCharTok{*}\NormalTok{alpha\_res}\SpecialCharTok{$}\NormalTok{total}\SpecialCharTok{$}\NormalTok{ase, }\DecValTok{3}\NormalTok{), }\StringTok{"]}\SpecialCharTok{\textbackslash{}n\textbackslash{}n}\StringTok{"}\NormalTok{)}
\end{Highlighting}
\end{Shaded}

\begin{verbatim}
## 95% CI:                      [ -0.766 , -0.545 ]
\end{verbatim}

\begin{Shaded}
\begin{Highlighting}[]
\FunctionTok{cat}\NormalTok{(}\StringTok{"Item{-}level statistics (what each item contributes):}\SpecialCharTok{\textbackslash{}n}\StringTok{"}\NormalTok{)}
\end{Highlighting}
\end{Shaded}

\begin{verbatim}
## Item-level statistics (what each item contributes):
\end{verbatim}

\begin{Shaded}
\begin{Highlighting}[]
\NormalTok{alpha\_res}\SpecialCharTok{$}\NormalTok{item.stats }\SpecialCharTok{|\textgreater{}}
\NormalTok{  dplyr}\SpecialCharTok{::}\FunctionTok{select}\NormalTok{(}\AttributeTok{r=}\NormalTok{r.cor, }\AttributeTok{r.drop=}\NormalTok{r.drop) }\SpecialCharTok{|\textgreater{}}
  \FunctionTok{round}\NormalTok{(}\DecValTok{3}\NormalTok{) }\SpecialCharTok{|\textgreater{}}
  \FunctionTok{print}\NormalTok{()}
\end{Highlighting}
\end{Shaded}

\begin{verbatim}
##         r r.drop
## E1 -1.559 -0.269
## E2 -2.252 -0.369
## E3  1.909 -0.005
## E4  0.789 -0.208
## E5  1.307 -0.049
\end{verbatim}

\textbf{Reading the alpha item statistics:}

\begin{itemize}
\tightlist
\item
  \textbf{r (item-total correlation):} How strongly each item correlates
  with the total scale score. Items with r \textless{} .30 are weakly
  related to the scale; consider dropping them.
\item
  \textbf{r.drop (alpha if item dropped):} What alpha would be if this
  item were removed. If this value is \emph{larger} than the overall
  alpha, the item is hurting reliability --- it may be measuring
  something different.
\end{itemize}

\subsection{McDonald's Omega --- The Better
Coefficient}\label{mcdonalds-omega-the-better-coefficient}

McDonald's omega does not assume equal loadings. It uses the actual
factor structure to estimate reliability. \texttt{omega\_h} is the
reliability attributable to the common factor; \texttt{omega\_total} is
the total reliability.

\begin{Shaded}
\begin{Highlighting}[]
\NormalTok{omega\_res }\OtherTok{\textless{}{-}} \FunctionTok{omega}\NormalTok{(extra\_items, }\AttributeTok{nfactors=}\DecValTok{1}\NormalTok{,}
                   \AttributeTok{title=}\StringTok{"Omega: Extraversion"}\NormalTok{, }\AttributeTok{plot=}\ConstantTok{FALSE}\NormalTok{)}

\FunctionTok{cat}\NormalTok{(}\StringTok{"=== McDonald\textquotesingle{}s Omega: Extraversion ===}\SpecialCharTok{\textbackslash{}n\textbackslash{}n}\StringTok{"}\NormalTok{)}
\end{Highlighting}
\end{Shaded}

\begin{verbatim}
## === McDonald's Omega: Extraversion ===
\end{verbatim}

\begin{Shaded}
\begin{Highlighting}[]
\FunctionTok{cat}\NormalTok{(}\StringTok{"omega\_h (hierarchical):  "}\NormalTok{, }\FunctionTok{round}\NormalTok{(omega\_res}\SpecialCharTok{$}\NormalTok{omega\_h, }\DecValTok{3}\NormalTok{), }\StringTok{"}\SpecialCharTok{\textbackslash{}n}\StringTok{"}\NormalTok{)}
\end{Highlighting}
\end{Shaded}

\begin{verbatim}
## omega_h (hierarchical):   0.767
\end{verbatim}

\begin{Shaded}
\begin{Highlighting}[]
\FunctionTok{cat}\NormalTok{(}\StringTok{"omega\_total:             "}\NormalTok{, }\FunctionTok{round}\NormalTok{(omega\_res}\SpecialCharTok{$}\NormalTok{omega.tot, }\DecValTok{3}\NormalTok{), }\StringTok{"}\SpecialCharTok{\textbackslash{}n}\StringTok{"}\NormalTok{)}
\end{Highlighting}
\end{Shaded}

\begin{verbatim}
## omega_total:              0.768
\end{verbatim}

\begin{Shaded}
\begin{Highlighting}[]
\FunctionTok{cat}\NormalTok{(}\StringTok{"Cronbach\textquotesingle{}s alpha:        "}\NormalTok{, }\FunctionTok{round}\NormalTok{(alpha\_res}\SpecialCharTok{$}\NormalTok{total}\SpecialCharTok{$}\NormalTok{raw\_alpha, }\DecValTok{3}\NormalTok{), }\StringTok{"}\SpecialCharTok{\textbackslash{}n\textbackslash{}n}\StringTok{"}\NormalTok{)}
\end{Highlighting}
\end{Shaded}

\begin{verbatim}
## Cronbach's alpha:         -0.655
\end{verbatim}

\begin{Shaded}
\begin{Highlighting}[]
\FunctionTok{cat}\NormalTok{(}\StringTok{"Interpretation:}\SpecialCharTok{\textbackslash{}n}\StringTok{"}\NormalTok{)}
\end{Highlighting}
\end{Shaded}

\begin{verbatim}
## Interpretation:
\end{verbatim}

\begin{Shaded}
\begin{Highlighting}[]
\FunctionTok{cat}\NormalTok{(}\StringTok{"  omega\_total and alpha are similar when loadings are equal (tau{-}equivalence holds)}\SpecialCharTok{\textbackslash{}n}\StringTok{"}\NormalTok{)}
\end{Highlighting}
\end{Shaded}

\begin{verbatim}
##   omega_total and alpha are similar when loadings are equal (tau-equivalence holds)
\end{verbatim}

\begin{Shaded}
\begin{Highlighting}[]
\FunctionTok{cat}\NormalTok{(}\StringTok{"  omega\_total \textgreater{} alpha when some items load higher than others (common situation)}\SpecialCharTok{\textbackslash{}n}\StringTok{"}\NormalTok{)}
\end{Highlighting}
\end{Shaded}

\begin{verbatim}
##   omega_total > alpha when some items load higher than others (common situation)
\end{verbatim}

\begin{Shaded}
\begin{Highlighting}[]
\FunctionTok{cat}\NormalTok{(}\StringTok{"  omega\_h is more conservative — it attributes only the common factor variance}\SpecialCharTok{\textbackslash{}n}\StringTok{"}\NormalTok{)}
\end{Highlighting}
\end{Shaded}

\begin{verbatim}
##   omega_h is more conservative — it attributes only the common factor variance
\end{verbatim}

\begin{Shaded}
\begin{Highlighting}[]
\FunctionTok{cat}\NormalTok{(}\StringTok{"  Best practice: report both alpha (for comparability) and omega (for accuracy)}\SpecialCharTok{\textbackslash{}n}\StringTok{"}\NormalTok{)}
\end{Highlighting}
\end{Shaded}

\begin{verbatim}
##   Best practice: report both alpha (for comparability) and omega (for accuracy)
\end{verbatim}

\begin{center}\rule{0.5\linewidth}{0.5pt}\end{center}

\section{Checking Assumptions and Diagnosing
Problems}\label{checking-assumptions-and-diagnosing-problems}

\subsection{Modification Indices}\label{modification-indices}

Modification indices (MIs) estimate how much chi-squared would drop if a
specific parameter were freed. Large MIs suggest parts of the model that
fit poorly.

\begin{Shaded}
\begin{Highlighting}[]
\NormalTok{mi }\OtherTok{\textless{}{-}} \FunctionTok{modindices}\NormalTok{(cfa\_fit, }\AttributeTok{sort.=}\ConstantTok{TRUE}\NormalTok{, }\AttributeTok{maximum.number=}\DecValTok{12}\NormalTok{)}

\NormalTok{mi }\SpecialCharTok{|\textgreater{}}
  \FunctionTok{filter}\NormalTok{(mi }\SpecialCharTok{\textgreater{}} \DecValTok{8}\NormalTok{) }\SpecialCharTok{|\textgreater{}}
\NormalTok{  dplyr}\SpecialCharTok{::}\FunctionTok{select}\NormalTok{(lhs, op, rhs, mi, epc, sepc.all) }\SpecialCharTok{|\textgreater{}}
  \FunctionTok{mutate}\NormalTok{(}
    \AttributeTok{mi       =} \FunctionTok{round}\NormalTok{(mi, }\DecValTok{2}\NormalTok{),}
    \AttributeTok{epc      =} \FunctionTok{round}\NormalTok{(epc, }\DecValTok{3}\NormalTok{),}
    \AttributeTok{sepc.all =} \FunctionTok{round}\NormalTok{(sepc.all, }\DecValTok{3}\NormalTok{),}
    \AttributeTok{Type     =} \FunctionTok{case\_when}\NormalTok{(}
\NormalTok{      op }\SpecialCharTok{==} \StringTok{"=\textasciitilde{}"} \SpecialCharTok{\textasciitilde{}} \StringTok{"Cross{-}loading: item loads on additional factor"}\NormalTok{,}
\NormalTok{      op }\SpecialCharTok{==} \StringTok{"\textasciitilde{}\textasciitilde{}"} \SpecialCharTok{\&}\NormalTok{ lhs }\SpecialCharTok{!=}\NormalTok{ rhs }\SpecialCharTok{\textasciitilde{}} \StringTok{"Residual correlation: items share specific variance"}\NormalTok{,}
\NormalTok{      op }\SpecialCharTok{==} \StringTok{"\textasciitilde{}\textasciitilde{}"} \SpecialCharTok{\&}\NormalTok{ lhs }\SpecialCharTok{==}\NormalTok{ rhs }\SpecialCharTok{\textasciitilde{}} \StringTok{"Free residual variance"}\NormalTok{,}
      \ConstantTok{TRUE} \SpecialCharTok{\textasciitilde{}} \StringTok{""}
\NormalTok{    )}
\NormalTok{  ) }\SpecialCharTok{|\textgreater{}}
  \FunctionTok{kbl}\NormalTok{(}\AttributeTok{caption=}\StringTok{"Modification Indices \textgreater{} 8 (Top Misspecifications)"}\NormalTok{,}
      \AttributeTok{col.names=}\FunctionTok{c}\NormalTok{(}\StringTok{"LHS"}\NormalTok{,}\StringTok{"Op"}\NormalTok{,}\StringTok{"RHS"}\NormalTok{,}\StringTok{"MI"}\NormalTok{,}\StringTok{"EPC"}\NormalTok{,}\StringTok{"Std.EPC"}\NormalTok{,}\StringTok{"Type"}\NormalTok{)) }\SpecialCharTok{|\textgreater{}}
  \FunctionTok{kable\_styling}\NormalTok{(}\AttributeTok{bootstrap\_options=}\FunctionTok{c}\NormalTok{(}\StringTok{"striped"}\NormalTok{,}\StringTok{"hover"}\NormalTok{,}\StringTok{"condensed"}\NormalTok{)) }\SpecialCharTok{|\textgreater{}}
  \FunctionTok{footnote}\NormalTok{(}\AttributeTok{general=}\FunctionTok{paste}\NormalTok{(}
    \StringTok{"MI = expected improvement in chi{-}squared if parameter is freed."}\NormalTok{,}
    \StringTok{"EPC = expected parameter change (how large the new parameter would be)."}\NormalTok{,}
    \StringTok{"}\SpecialCharTok{\textbackslash{}n}\StringTok{Only add parameters with THEORETICAL justification."}\NormalTok{,}
    \StringTok{"Never add parameters just because MI is large."}
\NormalTok{  ))}
\end{Highlighting}
\end{Shaded}

\begin{table}
\centering
\caption{\label{tab:mod-indices}Modification Indices > 8 (Top Misspecifications)}
\centering
\begin{tabular}[t]{l|l|l|r|r|r|l}
\hline
LHS & Op & RHS & MI & EPC & Std.EPC & Type\\
\hline
N1 & \textasciitilde{}\textasciitilde{} & N2 & 418.82 & 0.841 & 1.033 & Residual correlation: items share specific variance\\
\hline
Extraversion & =\textasciitilde{} & N4 & 200.79 & 0.448 & 0.285 & Cross-loading: item loads on additional factor\\
\hline
Openness & =\textasciitilde{} & E3 & 153.71 & 0.427 & 0.316 & Cross-loading: item loads on additional factor\\
\hline
N3 & \textasciitilde{}\textasciitilde{} & N4 & 134.10 & 0.403 & 0.283 & Residual correlation: items share specific variance\\
\hline
Openness & =\textasciitilde{} & E4 & 122.56 & -0.404 & -0.276 & Cross-loading: item loads on additional factor\\
\hline
Conscientiousness & =\textasciitilde{} & E5 & 121.50 & 0.343 & 0.255 & Cross-loading: item loads on additional factor\\
\hline
Extraversion & =\textasciitilde{} & O3 & 114.20 & -0.395 & -0.328 & Cross-loading: item loads on additional factor\\
\hline
Extraversion & =\textasciitilde{} & O4 & 113.87 & 0.343 & 0.287 & Cross-loading: item loads on additional factor\\
\hline
Neuroticism & =\textasciitilde{} & C5 & 108.75 & 0.352 & 0.216 & Cross-loading: item loads on additional factor\\
\hline
Extraversion & =\textasciitilde{} & A5 & 108.59 & -0.449 & -0.354 & Cross-loading: item loads on additional factor\\
\hline
Neuroticism & =\textasciitilde{} & C2 & 107.01 & 0.285 & 0.216 & Cross-loading: item loads on additional factor\\
\hline
C1 & \textasciitilde{}\textasciitilde{} & C2 & 106.98 & 0.288 & 0.263 & Residual correlation: items share specific variance\\
\hline
\multicolumn{7}{l}{\rule{0pt}{1em}\textit{Note: }}\\
\multicolumn{7}{l}{\rule{0pt}{1em}makecell[l]{MI = expected improvement in chi-squared if parameter is freed. EPC = expected parameter change (how large the new parameter would be). \Only add parameters with THEORETICAL justification. Never add parameters just because MI is large.}}\\
\end{tabular}
\end{table}

\textbf{What to do with modification indices:}

\begin{enumerate}
\def\labelenumi{\arabic{enumi}.}
\tightlist
\item
  Find the largest MI.
\item
  Ask: is there a \emph{theoretical reason} why this parameter should
  exist?
\item
  If yes, add it to the model and refit. Note it in your methods as a
  ``post-hoc modification.''
\item
  If no, do not add it. Large MI with no theoretical basis =
  capitalising on sample-specific noise.
\end{enumerate}

The most common justified modification is a \textbf{residual
correlation} between two items that share specific wording (e.g., both
contain the word ``feel'') or a specific subdimension beyond the common
factor.

\begin{center}\rule{0.5\linewidth}{0.5pt}\end{center}

\section{Reporting Standards}\label{reporting-standards}

When you report EFA or CFA results in a paper, include:

\begin{table}
\centering
\caption{\label{tab:reporting-template}Reporting Standards for EFA and CFA}
\centering
\begin{tabu} to \linewidth {>{\raggedright}X>{\raggedright}X}
\hline
What to Report & Example Language\\
\hline
\textbf{EFA: Extraction method and rotation} & Minimum residual extraction with oblimin rotation\\
\hline
\textbf{EFA: Number of factors and retention criterion} & Five factors retained based on parallel analysis (PA suggested k=5)\\
\hline
\textbf{EFA: Full loading matrix (all factors × all items)} & See Table X [loadings < |.30| suppressed for readability]\\
\hline
\textbf{EFA: Factor correlations (if oblique)} & Factor correlations ranged from r = .08 to r = .43 (Table Y)\\
\hline
\textbf{EFA: Variance explained per factor and cumulative} & The five factors explained 11, 9, 8, 8, and 7\% of variance (43\% cumulative)\\
\hline
\textbf{CFA: Sample size and missing data handling} & N = 2,236 complete cases (listwise deletion of 564 cases, 20.1\%)\\
\hline
\textbf{CFA: Estimator used} & MLR (Maximum Likelihood with robust standard errors)\\
\hline
\textbf{CFA: Model fit indices (CFI, TLI, RMSEA + CI, SRMR)} & CFI = .93, TLI = .92, RMSEA = .053 [.048, .058], SRMR = .063\\
\hline
\textbf{CFA: All standardised loadings with SE or CI} & Standardised loadings ranged from .32 to .81 (Table X)\\
\hline
\textbf{CFA: Factor correlations} & Factor correlations ranged from −.28 to .41 (Table Y)\\
\hline
\textbf{CFA: Any post-hoc modifications and justification} & Residual correlation allowed between E1 and E2 (shared wording; MI = 28.4)\\
\hline
\textbf{Reliability: Alpha and omega for each subscale} & Extraversion: α = .76, ω = .78; Agreeableness: α = .70, ω = .72\\
\hline
\end{tabu}
\end{table}

\begin{center}\rule{0.5\linewidth}{0.5pt}\end{center}

\section{A Complete Single-Scale
Example}\label{a-complete-single-scale-example}

To cement the workflow, here is a full single-scale analysis from start
to finish using only the Neuroticism items:

\begin{Shaded}
\begin{Highlighting}[]
\CommentTok{\# The Neuroticism subscale (N1–N5)}
\NormalTok{neuro }\OtherTok{\textless{}{-}}\NormalTok{ bfi\_complete[, }\FunctionTok{paste0}\NormalTok{(}\StringTok{"N"}\NormalTok{, }\DecValTok{1}\SpecialCharTok{:}\DecValTok{5}\NormalTok{)]}

\FunctionTok{cat}\NormalTok{(}\StringTok{"=== Complete Scale Analysis: Neuroticism (N1–N5) ===}\SpecialCharTok{\textbackslash{}n\textbackslash{}n}\StringTok{"}\NormalTok{)}
\end{Highlighting}
\end{Shaded}

\begin{verbatim}
## === Complete Scale Analysis: Neuroticism (N1–N5) ===
\end{verbatim}

\begin{Shaded}
\begin{Highlighting}[]
\CommentTok{\# 1. Describe}
\FunctionTok{cat}\NormalTok{(}\StringTok{"1. Descriptive statistics:}\SpecialCharTok{\textbackslash{}n}\StringTok{"}\NormalTok{)}
\end{Highlighting}
\end{Shaded}

\begin{verbatim}
## 1. Descriptive statistics:
\end{verbatim}

\begin{Shaded}
\begin{Highlighting}[]
\FunctionTok{describe}\NormalTok{(neuro)[, }\FunctionTok{c}\NormalTok{(}\StringTok{"n"}\NormalTok{,}\StringTok{"mean"}\NormalTok{,}\StringTok{"sd"}\NormalTok{,}\StringTok{"skew"}\NormalTok{,}\StringTok{"kurtosis"}\NormalTok{,}\StringTok{"min"}\NormalTok{,}\StringTok{"max"}\NormalTok{)] }\SpecialCharTok{|\textgreater{}}
  \FunctionTok{round}\NormalTok{(}\DecValTok{2}\NormalTok{) }\SpecialCharTok{|\textgreater{}}
  \FunctionTok{print}\NormalTok{()}
\end{Highlighting}
\end{Shaded}

\begin{verbatim}
##       n mean   sd  skew kurtosis min max
## N1 2436 2.94 1.58  0.37    -1.02   1   6
## N2 2436 3.52 1.53 -0.08    -1.06   1   6
## N3 2436 3.22 1.59  0.14    -1.18   1   6
## N4 2436 3.20 1.57  0.20    -1.09   1   6
## N5 2436 2.97 1.62  0.38    -1.07   1   6
\end{verbatim}

\begin{Shaded}
\begin{Highlighting}[]
\CommentTok{\# 2. Correlation matrix}
\FunctionTok{cat}\NormalTok{(}\StringTok{"}\SpecialCharTok{\textbackslash{}n}\StringTok{2. Item intercorrelation matrix:}\SpecialCharTok{\textbackslash{}n}\StringTok{"}\NormalTok{)}
\end{Highlighting}
\end{Shaded}

\begin{verbatim}
## 
## 2. Item intercorrelation matrix:
\end{verbatim}

\begin{Shaded}
\begin{Highlighting}[]
\FunctionTok{round}\NormalTok{(}\FunctionTok{cor}\NormalTok{(neuro), }\DecValTok{2}\NormalTok{) }\SpecialCharTok{|\textgreater{}} \FunctionTok{print}\NormalTok{()}
\end{Highlighting}
\end{Shaded}

\begin{verbatim}
##      N1   N2   N3   N4   N5
## N1 1.00 0.72 0.57 0.41 0.38
## N2 0.72 1.00 0.55 0.39 0.35
## N3 0.57 0.55 1.00 0.52 0.43
## N4 0.41 0.39 0.52 1.00 0.40
## N5 0.38 0.35 0.43 0.40 1.00
\end{verbatim}

\begin{Shaded}
\begin{Highlighting}[]
\FunctionTok{cat}\NormalTok{(}\StringTok{"   → All positive: items are in the same direction (no recoding issues)}\SpecialCharTok{\textbackslash{}n}\StringTok{"}\NormalTok{)}
\end{Highlighting}
\end{Shaded}

\begin{verbatim}
##    → All positive: items are in the same direction (no recoding issues)
\end{verbatim}

\begin{Shaded}
\begin{Highlighting}[]
\FunctionTok{cat}\NormalTok{(}\StringTok{"   → Correlations in range .25–.60: reasonable for a personality scale}\SpecialCharTok{\textbackslash{}n}\StringTok{"}\NormalTok{)}
\end{Highlighting}
\end{Shaded}

\begin{verbatim}
##    → Correlations in range .25–.60: reasonable for a personality scale
\end{verbatim}

\begin{Shaded}
\begin{Highlighting}[]
\CommentTok{\# 3. KMO}
\FunctionTok{cat}\NormalTok{(}\StringTok{"}\SpecialCharTok{\textbackslash{}n}\StringTok{3. KMO:"}\NormalTok{, }\FunctionTok{round}\NormalTok{(}\FunctionTok{KMO}\NormalTok{(neuro)}\SpecialCharTok{$}\NormalTok{MSA, }\DecValTok{3}\NormalTok{), }\StringTok{"}\SpecialCharTok{\textbackslash{}n}\StringTok{"}\NormalTok{)}
\end{Highlighting}
\end{Shaded}

\begin{verbatim}
## 
## 3. KMO: 0.795
\end{verbatim}

\begin{Shaded}
\begin{Highlighting}[]
\CommentTok{\# 4. EFA (1 factor)}
\FunctionTok{cat}\NormalTok{(}\StringTok{"}\SpecialCharTok{\textbackslash{}n}\StringTok{4. EFA — 1 factor:}\SpecialCharTok{\textbackslash{}n}\StringTok{"}\NormalTok{)}
\end{Highlighting}
\end{Shaded}

\begin{verbatim}
## 
## 4. EFA — 1 factor:
\end{verbatim}

\begin{Shaded}
\begin{Highlighting}[]
\NormalTok{efa1 }\OtherTok{\textless{}{-}} \FunctionTok{fa}\NormalTok{(neuro, }\AttributeTok{nfactors=}\DecValTok{1}\NormalTok{, }\AttributeTok{rotate=}\StringTok{"none"}\NormalTok{, }\AttributeTok{fm=}\StringTok{"minres"}\NormalTok{)}
\FunctionTok{cat}\NormalTok{(}\StringTok{"   Loadings:}\SpecialCharTok{\textbackslash{}n}\StringTok{"}\NormalTok{)}
\end{Highlighting}
\end{Shaded}

\begin{verbatim}
##    Loadings:
\end{verbatim}

\begin{Shaded}
\begin{Highlighting}[]
\FunctionTok{round}\NormalTok{(}\FunctionTok{unclass}\NormalTok{(efa1}\SpecialCharTok{$}\NormalTok{loadings), }\DecValTok{3}\NormalTok{) }\SpecialCharTok{|\textgreater{}} \FunctionTok{print}\NormalTok{()}
\end{Highlighting}
\end{Shaded}

\begin{verbatim}
##      MR1
## N1 0.792
## N2 0.764
## N3 0.763
## N4 0.598
## N5 0.530
\end{verbatim}

\begin{Shaded}
\begin{Highlighting}[]
\FunctionTok{cat}\NormalTok{(}\StringTok{"   Communalities (h2):}\SpecialCharTok{\textbackslash{}n}\StringTok{"}\NormalTok{)}
\end{Highlighting}
\end{Shaded}

\begin{verbatim}
##    Communalities (h2):
\end{verbatim}

\begin{Shaded}
\begin{Highlighting}[]
\FunctionTok{round}\NormalTok{(efa1}\SpecialCharTok{$}\NormalTok{communality, }\DecValTok{3}\NormalTok{) }\SpecialCharTok{|\textgreater{}} \FunctionTok{print}\NormalTok{()}
\end{Highlighting}
\end{Shaded}

\begin{verbatim}
##    N1    N2    N3    N4    N5 
## 0.627 0.583 0.582 0.357 0.281
\end{verbatim}

\begin{Shaded}
\begin{Highlighting}[]
\CommentTok{\# 5. CFA}
\FunctionTok{cat}\NormalTok{(}\StringTok{"}\SpecialCharTok{\textbackslash{}n}\StringTok{5. CFA — fit indices:}\SpecialCharTok{\textbackslash{}n}\StringTok{"}\NormalTok{)}
\end{Highlighting}
\end{Shaded}

\begin{verbatim}
## 
## 5. CFA — fit indices:
\end{verbatim}

\begin{Shaded}
\begin{Highlighting}[]
\NormalTok{neuro\_cfa\_model }\OtherTok{\textless{}{-}} \StringTok{\textquotesingle{}Neuroticism =\textasciitilde{} N1 + N2 + N3 + N4 + N5\textquotesingle{}}
\NormalTok{neuro\_cfa }\OtherTok{\textless{}{-}} \FunctionTok{cfa}\NormalTok{(neuro\_cfa\_model, }\AttributeTok{data=}\NormalTok{neuro, }\AttributeTok{std.lv=}\ConstantTok{TRUE}\NormalTok{)}
\NormalTok{fm }\OtherTok{\textless{}{-}} \FunctionTok{fitMeasures}\NormalTok{(neuro\_cfa, }\FunctionTok{c}\NormalTok{(}\StringTok{"cfi"}\NormalTok{,}\StringTok{"tli"}\NormalTok{,}\StringTok{"rmsea"}\NormalTok{,}\StringTok{"srmr"}\NormalTok{)) }\SpecialCharTok{|\textgreater{}} \FunctionTok{round}\NormalTok{(}\DecValTok{3}\NormalTok{)}
\FunctionTok{cat}\NormalTok{(}\StringTok{"   CFI ="}\NormalTok{, fm[}\StringTok{"cfi"}\NormalTok{], }\StringTok{"| TLI ="}\NormalTok{, fm[}\StringTok{"tli"}\NormalTok{],}
    \StringTok{"| RMSEA ="}\NormalTok{, fm[}\StringTok{"rmsea"}\NormalTok{], }\StringTok{"| SRMR ="}\NormalTok{, fm[}\StringTok{"srmr"}\NormalTok{], }\StringTok{"}\SpecialCharTok{\textbackslash{}n}\StringTok{"}\NormalTok{)}
\end{Highlighting}
\end{Shaded}

\begin{verbatim}
##    CFI = 0.923 | TLI = 0.847 | RMSEA = 0.167 | SRMR = 0.058
\end{verbatim}

\begin{Shaded}
\begin{Highlighting}[]
\CommentTok{\# 6. Reliability}
\FunctionTok{cat}\NormalTok{(}\StringTok{"}\SpecialCharTok{\textbackslash{}n}\StringTok{6. Reliability:}\SpecialCharTok{\textbackslash{}n}\StringTok{"}\NormalTok{)}
\end{Highlighting}
\end{Shaded}

\begin{verbatim}
## 
## 6. Reliability:
\end{verbatim}

\begin{Shaded}
\begin{Highlighting}[]
\FunctionTok{cat}\NormalTok{(}\StringTok{"   Cronbach\textquotesingle{}s alpha:"}\NormalTok{, }\FunctionTok{round}\NormalTok{(}\FunctionTok{alpha}\NormalTok{(neuro)}\SpecialCharTok{$}\NormalTok{total}\SpecialCharTok{$}\NormalTok{raw\_alpha, }\DecValTok{3}\NormalTok{), }\StringTok{"}\SpecialCharTok{\textbackslash{}n}\StringTok{"}\NormalTok{)}
\end{Highlighting}
\end{Shaded}

\begin{verbatim}
##    Cronbach's alpha: 0.817
\end{verbatim}

\begin{Shaded}
\begin{Highlighting}[]
\FunctionTok{cat}\NormalTok{(}\StringTok{"   McDonald\textquotesingle{}s omega:"}\NormalTok{, }\FunctionTok{round}\NormalTok{(}\FunctionTok{omega}\NormalTok{(neuro, }\AttributeTok{plot=}\ConstantTok{FALSE}\NormalTok{)}\SpecialCharTok{$}\NormalTok{omega.tot, }\DecValTok{3}\NormalTok{), }\StringTok{"}\SpecialCharTok{\textbackslash{}n}\StringTok{"}\NormalTok{)}
\end{Highlighting}
\end{Shaded}

\begin{verbatim}
##    McDonald's omega: 0.855
\end{verbatim}

\begin{Shaded}
\begin{Highlighting}[]
\FunctionTok{cat}\NormalTok{(}\StringTok{"}\SpecialCharTok{\textbackslash{}n}\StringTok{=== Summary ===}\SpecialCharTok{\textbackslash{}n}\StringTok{"}\NormalTok{)}
\end{Highlighting}
\end{Shaded}

\begin{verbatim}
## 
## === Summary ===
\end{verbatim}

\begin{Shaded}
\begin{Highlighting}[]
\FunctionTok{cat}\NormalTok{(}\StringTok{"Neuroticism items form a reliable, well{-}fitting single{-}factor scale.}\SpecialCharTok{\textbackslash{}n}\StringTok{"}\NormalTok{)}
\end{Highlighting}
\end{Shaded}

\begin{verbatim}
## Neuroticism items form a reliable, well-fitting single-factor scale.
\end{verbatim}

\begin{Shaded}
\begin{Highlighting}[]
\FunctionTok{cat}\NormalTok{(}\StringTok{"All items have acceptable loadings, communalities, and item{-}total correlations.}\SpecialCharTok{\textbackslash{}n}\StringTok{"}\NormalTok{)}
\end{Highlighting}
\end{Shaded}

\begin{verbatim}
## All items have acceptable loadings, communalities, and item-total correlations.
\end{verbatim}

\begin{center}\rule{0.5\linewidth}{0.5pt}\end{center}

\section{Common Mistakes}\label{common-mistakes}

{\def\LTcaptype{none} % do not increment counter
\begin{longtable}[]{@{}
  >{\raggedright\arraybackslash}p{(\linewidth - 4\tabcolsep) * \real{0.2647}}
  >{\raggedright\arraybackslash}p{(\linewidth - 4\tabcolsep) * \real{0.3824}}
  >{\raggedright\arraybackslash}p{(\linewidth - 4\tabcolsep) * \real{0.3529}}@{}}
\toprule\noalign{}
\begin{minipage}[b]{\linewidth}\raggedright
Mistake
\end{minipage} & \begin{minipage}[b]{\linewidth}\raggedright
Consequence
\end{minipage} & \begin{minipage}[b]{\linewidth}\raggedright
Prevention
\end{minipage} \\
\midrule\noalign{}
\endhead
\bottomrule\noalign{}
\endlastfoot
Confirming EFA structure on the same data & Circular --- CFA will fit
perfectly by design & Split-sample or hold-out validation \\
Using eigenvalue \textgreater{} 1 rule for retention & Typically
over-extracts factors & Use parallel analysis \\
Forcing orthogonal rotation when factors correlate & Distorted loadings;
factors misattributed & Check factor correlations; if \textgreater{} \\
Reporting only Cronbach's alpha & Misleads when loadings are unequal &
Report omega\_total alongside alpha \\
Adding MI-suggested parameters without theory & Overfitting; results
won't replicate & Only add theoretically justified parameters \\
Ignoring low communalities (h² \textless{} .20) & Weak items inflate
model misfit & Drop or revise items with h² \textless{} .20 \\
\end{longtable}
}

\begin{center}\rule{0.5\linewidth}{0.5pt}\end{center}

\section{Check Your Understanding}\label{check-your-understanding}

\textbf{1.} Run EFA on only the Agreeableness items (A1--A5). You should
find a one-factor solution. Report: (a) the KMO, (b) all five loadings,
(c) all five communalities, (d) the proportion of variance explained.
Interpret each.

\textbf{2.} Fit a CFA with the Agreeableness items as a single factor.
Report the fit indices. Are they within acceptable range? Which item has
the lowest loading? What would you do about it?

\textbf{3.} The residual correlation between N1 (``Get irritated
easily'') and N2 (``Seldom feel blue'') has a modification index of
35.2. Would you add this residual correlation to your CFA model? Explain
your reasoning.

\textbf{4.} 🟥 An EFA of 20 social attitude items produces a 4-factor
solution with Proportion Var values of .14, .11, .09, and .08
(cumulative = .42). A 5-factor solution adds a fifth factor with
Proportion Var = .04 (cumulative = .46). The BIC-parallel analysis
suggests 4 factors but the scree plot shows an elbow at 5. How do you
decide? What additional evidence would help?

\textbf{5.} 🟥 Compute Cronbach's alpha and McDonald's omega for the
Conscientiousness items (C1--C5). Then run a CFA with all five items
loading on one factor and check the loadings. Are alpha and omega
similar or different? What does the pattern of loadings tell you about
why they differ?

\begin{center}\rule{0.5\linewidth}{0.5pt}\end{center}

\begin{Shaded}
\begin{Highlighting}[]
\FunctionTok{sessionInfo}\NormalTok{()}
\DocumentationTok{\#\# R version 4.5.3 (2026{-}03{-}11 ucrt)}
\DocumentationTok{\#\# Platform: x86\_64{-}w64{-}mingw32/x64}
\DocumentationTok{\#\# Running under: Windows 11 x64 (build 26200)}
\DocumentationTok{\#\# }
\DocumentationTok{\#\# Matrix products: default}
\DocumentationTok{\#\#   LAPACK version 3.12.1}
\DocumentationTok{\#\# }
\DocumentationTok{\#\# locale:}
\DocumentationTok{\#\# [1] LC\_COLLATE=English\_Israel.utf8  LC\_CTYPE=English\_Israel.utf8   }
\DocumentationTok{\#\# [3] LC\_MONETARY=English\_Israel.utf8 LC\_NUMERIC=C                   }
\DocumentationTok{\#\# [5] LC\_TIME=English\_Israel.utf8    }
\DocumentationTok{\#\# }
\DocumentationTok{\#\# time zone: Asia/Jerusalem}
\DocumentationTok{\#\# tzcode source: internal}
\DocumentationTok{\#\# }
\DocumentationTok{\#\# attached base packages:}
\DocumentationTok{\#\# [1] stats     graphics  grDevices utils     datasets  methods   base     }
\DocumentationTok{\#\# }
\DocumentationTok{\#\# other attached packages:}
\DocumentationTok{\#\#  [1] kableExtra\_1.4.0     lubridate\_1.9.5      forcats\_1.0.1       }
\DocumentationTok{\#\#  [4] stringr\_1.6.0        dplyr\_1.2.0          readr\_2.2.0         }
\DocumentationTok{\#\#  [7] tidyr\_1.3.2          tidyverse\_2.0.0      GPArotation\_2025.3{-}1}
\DocumentationTok{\#\# [10] semTools\_0.5{-}8       semPlot\_1.1.8        lavaan\_0.6{-}21       }
\DocumentationTok{\#\# [13] psych\_2.6.3          purrr\_1.2.1          tibble\_3.3.1        }
\DocumentationTok{\#\# [16] ggplot2\_4.0.2       }
\DocumentationTok{\#\# }
\DocumentationTok{\#\# loaded via a namespace (and not attached):}
\DocumentationTok{\#\#   [1] RColorBrewer\_1.1{-}3    rstudioapi\_0.18.0     jsonlite\_2.0.0       }
\DocumentationTok{\#\#   [4] magrittr\_2.0.4        TH.data\_1.1{-}5         estimability\_1.5.1   }
\DocumentationTok{\#\#   [7] farver\_2.1.2          nloptr\_2.2.1          rmarkdown\_2.31       }
\DocumentationTok{\#\#  [10] vctrs\_0.7.2           minqa\_1.2.8           ClaudeR\_0.2.0        }
\DocumentationTok{\#\#  [13] base64enc\_0.1{-}6       htmltools\_0.5.9       Formula\_1.2{-}5        }
\DocumentationTok{\#\#  [16] htmlwidgets\_1.6.4     plyr\_1.8.9            sandwich\_3.1{-}1       }
\DocumentationTok{\#\#  [19] emmeans\_2.0.2         zoo\_1.8{-}15            igraph\_2.2.2         }
\DocumentationTok{\#\#  [22] mime\_0.13             lifecycle\_1.0.5       pkgconfig\_2.0.3      }
\DocumentationTok{\#\#  [25] Matrix\_1.7{-}4          R6\_2.6.1              fastmap\_1.2.0        }
\DocumentationTok{\#\#  [28] rbibutils\_2.4.1       shiny\_1.13.0          digest\_0.6.39        }
\DocumentationTok{\#\#  [31] OpenMx\_2.22.11        fdrtool\_1.2.18        colorspace\_2.1{-}2     }
\DocumentationTok{\#\#  [34] rprojroot\_2.1.1       textshaping\_1.0.5     Hmisc\_5.2{-}5          }
\DocumentationTok{\#\#  [37] labeling\_0.4.3        timechange\_0.4.0      abind\_1.4{-}8          }
\DocumentationTok{\#\#  [40] compiler\_4.5.3        here\_1.0.2            withr\_3.0.2          }
\DocumentationTok{\#\#  [43] glasso\_1.11           htmlTable\_2.4.3       S7\_0.2.1             }
\DocumentationTok{\#\#  [46] backports\_1.5.0       carData\_3.0{-}6         MASS\_7.3{-}65          }
\DocumentationTok{\#\#  [49] corpcor\_1.6.10        gtools\_3.9.5          tools\_4.5.3          }
\DocumentationTok{\#\#  [52] pbivnorm\_0.6.0        foreign\_0.8{-}91        otel\_0.2.0           }
\DocumentationTok{\#\#  [55] zip\_2.3.3             httpuv\_1.6.17         nnet\_7.3{-}20          }
\DocumentationTok{\#\#  [58] glue\_1.8.0            quadprog\_1.5{-}8        nlme\_3.1{-}168         }
\DocumentationTok{\#\#  [61] promises\_1.5.0        lisrelToR\_0.3         grid\_4.5.3           }
\DocumentationTok{\#\#  [64] checkmate\_2.3.4       cluster\_2.1.8.2       reshape2\_1.4.5       }
\DocumentationTok{\#\#  [67] generics\_0.1.4        gtable\_0.3.6          tzdb\_0.5.0           }
\DocumentationTok{\#\#  [70] data.table\_1.18.2.1   hms\_1.1.4             xml2\_1.5.2           }
\DocumentationTok{\#\#  [73] sem\_3.1{-}16            pillar\_1.11.1         rockchalk\_1.8.157    }
\DocumentationTok{\#\#  [76] later\_1.4.8           splines\_4.5.3         lattice\_0.22{-}9       }
\DocumentationTok{\#\#  [79] survival\_3.8{-}6        kutils\_1.73           tidyselect\_1.2.1     }
\DocumentationTok{\#\#  [82] miniUI\_0.1.2          pbapply\_1.7{-}4         knitr\_1.51           }
\DocumentationTok{\#\#  [85] reformulas\_0.4.4      gridExtra\_2.3         svglite\_2.2.2        }
\DocumentationTok{\#\#  [88] stats4\_4.5.3          xfun\_0.57             qgraph\_1.9.8         }
\DocumentationTok{\#\#  [91] arm\_1.15{-}2            stringi\_1.8.7         yaml\_2.3.12          }
\DocumentationTok{\#\#  [94] boot\_1.3{-}32           evaluate\_1.0.5        codetools\_0.2{-}20     }
\DocumentationTok{\#\#  [97] mi\_1.2                cli\_3.6.5             RcppParallel\_5.1.11{-}2}
\DocumentationTok{\#\# [100] rpart\_4.1.24          xtable\_1.8{-}8          systemfonts\_1.3.2    }
\DocumentationTok{\#\# [103] Rdpack\_2.6.6          Rcpp\_1.1.1            coda\_0.19{-}4.1        }
\DocumentationTok{\#\# [106] png\_0.1{-}9             XML\_3.99{-}0.23         parallel\_4.5.3       }
\DocumentationTok{\#\# [109] jpeg\_0.1{-}11           lme4\_2.0{-}1            viridisLite\_0.4.3    }
\DocumentationTok{\#\# [112] mvtnorm\_1.3{-}6         scales\_1.4.0          openxlsx\_4.2.8.1     }
\DocumentationTok{\#\# [115] rlang\_1.1.7           multcomp\_1.4{-}30       mnormt\_2.1.2}
\end{Highlighting}
\end{Shaded}


\end{document}
